\section{Quadratic arithmetic contrasts}
\label{app:qc}

Convex quadratic programming over a rational polyhedron is the arithmetic
baseline for this paper. Its optimal value and a minimum-norm optimizer are
rational, and the classical ellipsoid analysis computes them exactly in
polynomial time~\cite{KozlovTarasovKhachiyan1980}. The main text shows how far
this picture breaks for quartic objectives: exact comparison becomes
$\PosSLP$-complete, and exact optimizers can have exponential degree or
exponentially long rational coordinates. Convex quadratic \emph{constraints}
lie between these two regimes. This appendix records what they do to the
arithmetic of values, points, and certificates. The results calibrate the
quartic theorems: some phenomena of the main text already occur for convex
quadratic constraints, and others do not.

The structural parameter throughout is the dimension $k$ of the span of the
constraint Hessians. Section~\ref{app:qc-degree} proves that the optimal value
and a canonical optimizer lie in one number field whose degree is at most
$\mathcal B(n,k)=\max_{0\le s\le\min(k,n)}2^s\binom ns$, with coefficient
heights $L^{O(k+1)}$, and that this degree bound is attained. The remaining
sections use these bounds and contrast them with lower bounds in specified
output formats. Table~\ref{tab:qc-summary} lists the results.

\begin{table}[t]
\centering
\small
\begin{tabularx}{\textwidth}{@{}>{\raggedright\arraybackslash}p{3.6cm}
  >{\raggedright\arraybackslash}X
  >{\raggedright\arraybackslash}p{2.3cm}@{}}
\toprule
Question & Answer & Result\\
\midrule
Degree and height of value and canonical optimizer
 & one field of degree at most $\mathcal B(n,k)$; heights $L^{O(k+1)}$;
   degree attained by strictly feasible rational instances
 & Thms.~\ref{thm:qc-degree}, \ref{thm:qc-height}, \ref{thm:qc-sharp}\\
Coefficients in a number field of degree $D$
 & height and separation bounds linear in $D$ and in the input height
 & Thm.~\ref{thm:qc-height}\\
Rational feasible points
 & always exist, are dense, and are computable when $k\le2$; an irrational
   singleton exists when $k=3$
 & Thm.~\ref{thm:qc-span-two}, Prop.~\ref{prop:qc-span-three}\\
Singleton spectrahedra
 & one scalar variable: exactly the totally real numbers, minimum size
   $2\deg\alpha$; corank one: exactly one real embedding
 & Thms.~\ref{thm:qc-pencil}, \ref{thm:qc-corank-one}\\
Expanded output length
 & short inputs force degree $n^{\Omega(k)}$ and dense coefficient lists of
   minimal polynomials of the value and of an input coordinate, also for a
   singleton cut out by ellipsoids
 & Thms.~\ref{thm:qc-blocks}, \ref{thm:qc-sparse},
   Cor.~\ref{cor:qc-feasible-output}\\
One convex quartic row
 & value degree $3^n$
 & Prop.~\ref{prop:qc-quartic}\\
Common field from approximations
 & primitive element and coordinate polynomials with polynomial overhead
 & Thm.~\ref{thm:qc-recovery}\\
Optimality and infeasibility certificates
 & one field, at most $k$ curved reductions; rational infeasibility
   aggregates of length $L^{O(k+1)}$, exponential in $k$ when necessary
 & Thms.~\ref{thm:qc-certificate}, \ref{thm:qc-rational-infeasibility}\\
Affine data in a number field
 & exact optimizer in the input field in polynomial time
 & Thm.~\ref{thm:qc-number-field-qp}\\
Nonconvex integer rows
 & exponentially long integer output already in dimension two
 & Prop.~\ref{prop:qc-pell}\\
\bottomrule
\end{tabularx}
\caption{Quadratic arithmetic contrasts proved in this appendix. Here $n$ is
the number of continuous variables, $k$ the dimension of the span of the
native constraint Hessians, and $L$ the total binary input length.}
\label{tab:qc-summary}
\end{table}

\subsection{Setting, output contracts, and heights}
\label{app:qc-setting}

\begin{definition}[Native PSD quadratic systems]
\label{def:qc-native}
A \emph{native PSD quadratic system} is
\begin{equation}
\label{eq:qc-system}
F=\{x\in\R^n: Ax\le b,\ g_i(x)\le0\ (1\le i\le m)\},\qquad
g_i(x)=\tfrac12x^{\mathsf T}Q_ix+a_i^{\mathsf T}x+c_i,
\end{equation}
where $Q_i$ is a symmetric positive semidefinite matrix for every $i$. An
affine equality is written as two opposite rows. The \emph{Hessian span} of
the system is $k=\dim\Span\{Q_1,\ldots,Q_m\}$. An objective
$f(x)=\tfrac12x^{\mathsf T}Q_0x+a_0^{\mathsf T}x+c_0$ with $Q_0\succeq0$ may
be added; its Hessian is not counted in $k$.
\end{definition}

Unless a field is named, all data are rational and explicitly encoded, every
matrix entry counted. The total binary length is $L\ge\max\{2,n,m\}$. Printed
input lengths in this appendix refer to this native matrix encoding. The
explicit format of Definition~\ref{def:models-polynomial} lists only nonzero
coefficients but attaches an exponent vector of $O(n)$ bits to each of the at
most $(n+1)(n+2)/2$ monomials of a quadratic. The two lengths are
polynomially related, so bounds polynomial in $L$, such as $L^{O(k+1)}$, hold
in either format; the output lower bounds below state the input length in
both formats. The span dimension $k$ is computed by rational rank. It
differs from the number of rows, from the rank of $\sum_iQ_i$, and from the
dimension of the span of the complete polynomials $g_i$: the rows
$\norm{x}^2+2x_j-1$, $1\le j\le n$, have $k=1$.

The word \emph{native} matters. Positive semidefiniteness is required of the
polynomials as written, not of the set $F$. Squaring the second-order cone
constraints $\norm{(1,1)}_2\le t$ and $\norm{(t,t)}_2\le2$ gives
$2-t^2\le0$ and $2t^2-4\le0$; the pair forces $t=\sqrt2$, while the first
squared row has a negative Hessian. The rational-point theorem of
Section~\ref{app:qc-rational} does not apply to such rows.

If $F\ne\varnothing$ and $f$ attains its infimum on $F$, the optimal set is
nonempty, closed and convex. Its unique point of least Euclidean norm in the
original coordinates is the \emph{canonical optimizer}, denoted $p$, and
$\theta=f(p)$ is the optimal value. For $f=0$, $p$ is the minimum-norm
feasible point. Other points of a positive-dimensional optimal face can be
transcendental and are not described by the results below.

\paragraph{Output contracts.}
The exact representations are those of
Definition~\ref{def:models-representations}. Two algebraic variants are kept
apart: a \emph{coordinatewise} representation gives each coordinate's
minimal polynomial and isolating interval, and a \emph{common-field}
representation gives one primitive irreducible integer polynomial $P$ of
degree $d$, a rational interval isolating one real root $\alpha$ of $P$, and
rational polynomials $b_1,\ldots,b_n$ of degree less than $d$ with
$x_j=b_j(\alpha)$. Coordinatewise output does not bound the degree of the
field generated by all coordinates: square roots of $n$ distinct primes have
degree two each and generate a field of degree $2^n$. The common-field format
lets a verifier evaluate any rational polynomial at the point by univariate
arithmetic modulo $P$.

Five quantities are kept separate: the bit-operation time of an algorithm;
the expanded length of an output in a named format; algebraic degrees; the
data and verification cost of a certificate; and the accuracy $2^{-q}$
requested from an approximation oracle. Two degrees must be distinguished.
The degree of one scalar bounds from below the length of a dense list of
that scalar's minimal polynomial. The degree of the field generated by all
coordinates bounds from below the length of a common-field representation,
but not of a coordinatewise one, whose separate minimal polynomials can be
short when the joint degree is large. Neither degree bounds implicit
descriptions. A circuit over the rationals computes only rational numbers, so
for irrational coordinates the compact alternatives are implicit systems,
towers of radicals, or circuits with root-selection gates. Lower bounds in
this appendix always name the format that they constrain.

\paragraph{Heights.}
For a number field $M$ and a place $u$ of $M$, let $\abs{\cdot}_u$ be the
absolute value extending the ordinary or $\ell$-adic absolute value of $\Q$,
and let $n_u=[M_u:\Q_u]$. Archimedean places carry the ordinary modulus, with
$n_u=2$ at complex places, and $\sum_{u\mid\infty}n_u=[M:\Q]$. The product
formula $\prod_u\abs{x}_u^{n_u}=1$ holds for $x\in M^\times$; equivalently
$\sum_un_u\log\abs{x}_u=0$. The absolute logarithmic Weil height
of $x\in M$ and the affine and projective heights of a vector $y\in M^r$ are
\[
\height(x)=\frac1{[M:\Q]}\sum_un_u\log\max(1,\abs{x}_u),\qquad
\height_{\rm aff}(y)=\frac1{[M:\Q]}\sum_un_u\log\max\bigl(1,\max_j\abs{y_j}_u\bigr),
\]
and, for $y\ne0$,
$\height_{\rm proj}(y)=[M:\Q]^{-1}\sum_un_u\log\max_j\abs{y_j}_u$. These
quantities do not depend on the field $M$ containing the
entries~\cite{BombieriGubler2006}.

\begin{lemma}[Height calculus]
\label{lem:qc-heights}
Let $M$ be a number field and $x,x_1,\ldots,x_r\in M$.
\begin{enumerate}[label=(\alph*),leftmargin=*]
\item $\height(x_1x_2)\le\height(x_1)+\height(x_2)$,
$\height(x^{-1})=\height(x)$ for $x\ne0$,
$\height(x_1+\cdots+x_r)\le\sum_j\height(x_j)+\log r$, and
$\height_{\rm proj}(y)\le\height_{\rm aff}(y)\le\sum_j\height(y_j)$. If
$y_i,y_j\ne0$ are entries of $y$, then $\height(y_i/y_j)\le\height_{\rm proj}(y)$.
\item For a rational number $a/b$ in lowest terms,
$\height(a/b)=\log\max(\abs a,\abs b)$.
\item If $M\subset\R$ through a fixed real embedding and $x\ne0$, then
$\abs{\log\abs x}\le[M:\Q]\height(x)$.
\item For an $r\times r$ matrix $G$ with entries in $M$,
$\height(\det G)\le\sum_{i,j}\height(G_{ij})+\log r!$.
\end{enumerate}
\end{lemma}

\begin{proof}
For (a), use
\[
\begin{gathered}
\max(1,\abs{x_1x_2}_u)\le\max(1,\abs{x_1}_u)\max(1,\abs{x_2}_u),\\
\sum_un_u\log\max(1,\abs{x}_u^{-1})=\sum_un_u\log\max(1,\abs{x}_u),
\end{gathered}
\]
the second by the product formula. At an archimedean place $\abs{\sum_jx_j}_u\le r\max_j\abs{x_j}_u$, at the other
places the maximum suffices, and the archimedean weights sum to $[M:\Q]$. The
two vector inequalities are immediate from the definitions. For the ratio,
$\height(y_i/y_j)$ is the projective height of $(y_j,y_i)$, and the
maximum over two entries is at most the maximum over all entries. For (b),
a prime $\ell$ dividing $b$ contributes $v_\ell(b)\log\ell$, the other primes
contribute nothing, and these contributions sum to $\log\abs b$; the
archimedean term adds $\log\max(1,\abs{a/b})$. For (c), the
real place contributes $n_u\log\max(1,\abs x)\le[M:\Q]\height(x)$; apply this
also to $x^{-1}$. For (d), expand the determinant. At an archimedean place,
$\abs{\det G}_u\le r!\prod_{i,j}\max(1,\abs{G_{ij}}_u)$; at the other places
the factor $r!$ is absent. Summing with the weights gives the claim.
\end{proof}

We also use the following elementary passage from an annihilating
polynomial to a minimal polynomial. If $\vartheta$ is a root of a nonzero
$R\in\Z[w]$ of degree $e$ with all coefficients bounded by $\exp(W)$, $W\ge0$,
then the primitive minimal polynomial $m_\vartheta\in\Z[w]$ of $\vartheta$
divides $R$ by Gauss's lemma. Its leading coefficient divides that of $R$, and
its roots are roots of $R$, hence bounded by $1+\exp(W)$ by Cauchy's bound.
Expanding $m_\vartheta$ over its roots gives
\begin{equation}
\label{eq:qc-minpoly-height}
\log\norm{m_\vartheta}_\infty\le W+\deg m_\vartheta\,\bigl(\log2+\log(1+\exp W)\bigr)
\le(\deg m_\vartheta+1)(W+2).
\end{equation}
Finally, for a polynomial $\sum_{i=0}^e\pi_iw^i$ with $\pi_e\ne0$, every
complex root $z$ satisfies Fujiwara's bound
$\abs z\le2\max_{i<e}\abs{\pi_i/\pi_e}^{1/(e-i)}$: if $\abs z$ were larger,
the terms $\abs{\pi_iz^i}$ with $i<e$ would sum to less than
$\abs{\pi_e}\abs z^e\sum_{j\ge1}2^{-j}$. The same bound holds at every
archimedean place, and at a nonarchimedean place the ultrametric inequality
gives $\abs z_u\le\max(1,\max_{i<e}\abs{\pi_i/\pi_e}_u)$.

\subsection{Degree, height, and separation at fixed Hessian span}
\label{app:qc-degree}

This subsection proves the bounds that the rest of the appendix uses. They
explain why convex quadratic constraints are arithmetically mild when their
Hessians span few directions, even with arbitrarily many variables and rows,
no bounds, and no constraint qualification. The proof restricts to the active
rows at the canonical optimizer, compresses the multipliers to at most $k$,
and passes the resulting nonsingular polynomial system through two ordered
limits. Two elimination arguments are needed: a root count gives the sharp
degree, and a finite multiplication-matrix construction gives the height.

\paragraph{Setting over a number field.}
Let $K\subset\R$ be a number field of degree $D$, with the embedding fixed;
$K=\Q$ is the main case. In this subsection the data of~\eqref{eq:qc-system}
and of the objective $f$ lie in $K$, and $Q_0,Q_1,\ldots,Q_m\succeq0$ in the
fixed embedding. The span dimension of the $Q_i$ over $K$ equals the real span
dimension $k$, because the rank of the matrix of Hessian entries does not
depend on the field. Assume that $F\ne\varnothing$ and that $f$ attains its
infimum on $F$, and let $p$ and $\theta$ be the canonical optimizer and the
optimal value. Let $S\ge2$ bound $n$, the number of rows, and the number of
listed input scalars, and let $B\ge1$ bound the absolute heights of the input
scalars. For $K=\Q$ one may take $S=L$ and $B=L$.

\begin{lemma}[Active restriction]
\label{lem:qc-restriction}
Let $I_a$ and $I_q$ be the affine rows and the quadratic rows that are active
at $p$. Choose $\mathcal J\subseteq I_q$ such that $\{Q_j:j\in\mathcal J\}$ is a
basis of $\Span\{Q_i:i\in I_q\}$, and write $Q_i=\sum_{j\in\mathcal J}
\gamma_{ij}Q_j$ with $\gamma_{ij}\in K$. Let $H$ be the affine space defined
by the active affine rows taken as equalities and by the equations
\begin{equation}
\label{eq:qc-differences}
g_i(x)-\sum_{j\in\mathcal J}\gamma_{ij}g_j(x)=0\qquad(i\in I_q),
\end{equation}
whose left sides are affine. Let $x=x_0+Vu$, $u\in\R^d$, be a parametrization
of $H$ with $x_0\in K^n$ and $V\in K^{n\times d}$ of full column rank, and put
$r_i(u)=g_i(x_0+Vu)$ for $i\in I_q$, $\tilde f(u)=f(x_0+Vu)$, and
$N(u)=\norm{x_0+Vu}^2$. Then the following hold.
\begin{enumerate}[label=(\alph*),leftmargin=*]
\item $p=x_0+Vu^*$ for a unique $u^*$, and $r_i(u^*)=0$ for $i\in I_q$.
\item The polynomials $r_i$, $i\in I_q$, span a space of dimension at most
$k$ over $K$.
\item On $C=\{u\in\R^d:r_i(u)\le0\ (i\in I_q)\}$ one has $\tilde f\ge\theta$,
and $u^*$ is the unique minimizer of $N$ on $\{u\in C:\tilde f(u)=\theta\}$.
\item There is an absolute constant $c_0$ such that every entry of $x_0$ and
$V$, and every coefficient of $r_i$, $\tilde f$ and $N$, has height at most
$S^{c_0}(B+1)$. These bounds do not depend on which rows are active.
\end{enumerate}
\end{lemma}

\begin{proof}
Every active row vanishes at $p$, so $p\in H$ and (a) holds. Each difference
in~\eqref{eq:qc-differences} is affine and vanishes on $H$, so it vanishes
identically after substitution. Hence $r_i=\sum_j\gamma_{ij}r_j$ as
polynomials, proving (b).

For (c), let $F_2=\{x\in H:g_i(x)\le0\ (i\in I_q)\}$, the image of $C$. Let
$y\in F_2$ and $z_t=p+t(y-p)$ for $0<t\le1$. The active affine rows hold on
$H$, the rows in $I_q$ hold at $z_t$ by convexity, and every inactive row is
strictly slack at $p$ and therefore at $z_t$ for small $t$. Thus $z_t\in F$
for small $t$. If $f(y)<\theta$, convexity gives $f(z_t)<\theta$, a
contradiction. If $f(y)=\theta$ and $y\ne p$ with $\norm y\le\norm p$, then
$z_t$ is optimal for small $t$, and strict convexity of the squared norm
gives $\norm{z_t}<\norm p$, again a contradiction.

For (d), the $\gamma_{ij}$ solve a $K$-linear system whose matrix consists of
Hessian entries. Cramer's rule and Lemma~\ref{lem:qc-heights}(d) bound their
heights by $2(S^2B+S\log S)$. The coefficients of~\eqref{eq:qc-differences}
are sums of at most $S+1$ products of these numbers with input scalars. The
vectors $x_0$ and the columns of $V$ are Cramer quotients of minors of order
at most $n\le S$ of the equation matrix of $H$, after choosing the free
coordinates among the original ones. The restricted coefficients are sums of
at most $S^2$ products of three such numbers. Each of these finitely many
steps multiplies the preceding height bound by a fixed polynomial in $S$ and
adds a term of order $S\log S$, by Lemma~\ref{lem:qc-heights}(a),(d).
\end{proof}

If $d=0$, then $p=x_0\in K^n$, and all statements below are immediate. Assume
$d\ge1$ from now on.

\begin{lemma}[Ordered regularization]
\label{lem:qc-ordered}
For $\varepsilon>0$, the function $f_\varepsilon=\tilde f+\varepsilon N$ has
a unique minimizer $u_\varepsilon$ on $C$, and $x_0+Vu_\varepsilon\to p$ as
$\varepsilon\downarrow0$. For fixed $\varepsilon>0$ and $\delta>0$,
$f_\varepsilon$ has a unique minimizer $u_{\varepsilon,\delta}$ on
$C_\delta=\{u:r_i(u)\le\delta\ (i\in I_q)\}$, and
$u_{\varepsilon,\delta}\to u_\varepsilon$ as $\delta\downarrow0$. The point
$u^*$ satisfies every inequality of $C_\delta$ strictly.
\end{lemma}

\begin{proof}
The Hessian of $f_\varepsilon$ is $\nabla^2\tilde f+2\varepsilon V^{\mathsf T}V$,
which is positive definite because $V$ has full column rank. Hence
$f_\varepsilon$ is strongly convex and has unique minimizers on the nonempty
closed convex sets $C$ and $C_\delta$. Comparison with $u^*$ and
Lemma~\ref{lem:qc-restriction}(c) give
\[
\theta\le\tilde f(u_\varepsilon),\qquad
\tilde f(u_\varepsilon)+\varepsilon N(u_\varepsilon)\le\theta+\varepsilon\norm p^2,
\]
so $N(u_\varepsilon)\le\norm p^2$ and $\tilde f(u_\varepsilon)\to\theta$.
Every accumulation point $\bar u$ of $u_\varepsilon$ lies in $C$, has
$\tilde f(\bar u)=\theta$ and $N(\bar u)\le\norm p^2$, so $\bar u=u^*$ by
Lemma~\ref{lem:qc-restriction}(c). For fixed $\varepsilon$, the points
$u_{\varepsilon,\delta}$ lie in the compact set
$\{u:f_\varepsilon(u)\le f_\varepsilon(u_\varepsilon)\}$, since
$u_\varepsilon\in C_\delta$. An accumulation point as $\delta\downarrow0$ lies
in $C$ and does not exceed the minimum value on $C$, so it equals
$u_\varepsilon$. The last assertion holds because $r_i(u^*)=0<\delta$.
\end{proof}

Only the ordered limits, first $\delta\downarrow0$ and then
$\varepsilon\downarrow0$, are used below. No convergence along a joint
schedule of the two parameters is asserted.

For a set $J\subseteq I_q$ with $\abs J=s$, consider the square system in
$(u,\lambda)\in\R^{d}\times\R^{s}$
\begin{equation}
\label{eq:qc-kkt}
\Phi_J(u,\lambda;\varepsilon,\delta)=
\Bigl(\nabla f_\varepsilon(u)+\sum_{j\in J}\lambda_j\nabla r_j(u),\
(r_j(u)-\delta)_{j\in J}\Bigr)=0 .
\end{equation}
Write $\tilde Q_i$, $\tilde a_i$ and $\tilde c_i$ for the Hessian, linear
coefficient vector and constant of $r_i$ for $i\in I_q$, and of $\tilde f$ for
$i=0$.

\begin{lemma}[Small regular support]
\label{lem:qc-kkt}
For $\varepsilon,\delta>0$ there is $J\subseteq I_q$ with
$s=\abs J\le\min(k,d)$ and $\lambda\in\R^J_{>0}$ such that
$(u_{\varepsilon,\delta},\lambda)$ solves~\eqref{eq:qc-kkt} and the Jacobian
of $\Phi_J$ with respect to $(u,\lambda)$ is nonsingular there. For every
sequence $\varepsilon_\ell\downarrow0$ there are sequences
$\delta_{\ell,\nu}\downarrow0$ and one set $J$ such that this holds for all
pairs $(\varepsilon_\ell,\delta_{\ell,\nu})$, after passing to a subsequence
in $\ell$.
\end{lemma}

\begin{proof}
By Lemma~\ref{lem:qc-ordered}, $u^*$ is a strict feasible point of the convex
problem $\min\{f_\varepsilon(u):u\in C_\delta\}$, so the convex
Karush--Kuhn--Tucker conditions hold at its minimizer: $-\nabla f_\varepsilon$
is a nonnegative combination of the gradients of the rows with
$r_i=\delta$~\cite{Rockafellar1970}. Choose such a representation with
minimal support $J$. Its gradients are linearly independent: a dependence
$\sum_j\eta_j\nabla r_j=0$ with some $\eta_j>0$ lets us replace $\lambda$ by
$\lambda-\tau\eta$ for the largest $\tau$ keeping $\lambda\ge0$, which
removes an index. By Lemma~\ref{lem:qc-restriction}(b), the gradients of the
$r_i$ at one point span at most $k$ dimensions, so $s\le\min(k,d)$.

The Jacobian of~\eqref{eq:qc-kkt} is
$\left(\begin{smallmatrix}M&G\\G^{\mathsf T}&0\end{smallmatrix}\right)$ with
$M=\tilde Q_0+2\varepsilon V^{\mathsf T}V+\sum_{j\in J}\lambda_j\tilde Q_j\succ0$
and $G=(\nabla r_j(u_{\varepsilon,\delta}))_{j\in J}$ of full column rank. If
$Mx+Gy=0$ and $G^{\mathsf T}x=0$, then $x^{\mathsf T}Mx=0$, so $x=0$ and then
$y=0$. For the last assertion, there are finitely many subsets $J$: choose one
that occurs along a sequence $\delta_{\ell,\nu}\downarrow0$ for each $\ell$,
and then pass to a subsequence of $\ell$ on which this choice is constant.
\end{proof}

The support is fixed before any output is chosen. The multipliers may be
unbounded along the sequences; no bound on them is used.

\paragraph{The degree bound.}
The next lemma turns a uniform count of nonsingular roots into one polynomial
identity that holds at every real nonsingular root, including roots at
exceptional parameter values. Counting the roots of individual rational
specializations would not suffice: rational numbers can converge to a
transcendental number.

\begin{lemma}[One identity at all regular roots]
\label{lem:qc-parameter}
Let $K\subset\R$ be a field, $t=(t_1,\ldots,t_r)$ and $y=(y_1,\ldots,y_N)$
indeterminates, and $\Phi_1,\ldots,\Phi_N,\Gamma\in K[t,y]$. Let $\mathbb K=K(t)$
and suppose that $\Phi(t,\cdot)$ has at most $\beta$ nonsingular common zeros
in $\overline{\mathbb K}^N$. If some real $(t^\circ,y^\circ)$ is a common zero
with $\det\partial_y\Phi(t^\circ,y^\circ)\ne0$, then there is a nonzero
$R\in K[t,w]$ with $\deg_wR\le\beta$ such that $R(t,\Gamma(t,y))=0$ at every
real common zero $(t,y)$ of $\Phi$ with $\det\partial_y\Phi(t,y)\ne0$.
\end{lemma}

\begin{proof}
Let $\Delta=\det\partial_y\Phi$ and
$E=\mathbb K[y,\Delta^{-1}]/(\Phi_1,\ldots,\Phi_N)$. After extension of
scalars to $\overline{\mathbb K}$, the maximal ideals of $E$ correspond to
the common zeros $z$ with $\Delta(z)\ne0$. At such a zero the differentials of
the $\Phi_i$ span the cotangent space, so the $\Phi_i$ generate the maximal
ideal of the local ring of $\overline{\mathbb K}^N$ at $z$, and the local ring of
$E\otimes\overline{\mathbb K}$ at $z$ is the field $\overline{\mathbb K}$. Thus
every point of the affine scheme defined by $E\otimes\overline{\mathbb K}$ is
an isolated reduced point. A Noetherian ring with this property has finitely
many points and is the product of its residue fields, so
$\dim_{\mathbb K}E$ equals the number of nonsingular common zeros, which is at
most $\beta$.

The algebra $E$ is not zero. Otherwise $\Delta^e\in(\Phi)$ in
$\mathbb K[y]$ for some $e$, and clearing denominators gives
$c(t)\Delta^e=\sum_ih_i\Phi_i$ with $0\ne c\in K[t]$ and $h_i\in K[t,y]$. By the
implicit function theorem there is a real analytic branch $y(t)$ of
nonsingular zeros on a neighborhood of $t^\circ$, so $c$ vanishes on an open
set, which is impossible.

Multiplication by the class of $\Gamma$ on the $\mathbb K$-vector space $E$
has a monic characteristic polynomial $\chi\in\mathbb K[w]$ of degree
$\dim E\le\beta$, and $\chi(\Gamma)=0$ in $E$ by the Cayley--Hamilton
theorem. Clearing denominators gives a nonzero $R\in K[t,w]$ with
$\deg_wR=\dim E$ and an identity
$c_2(t)\Delta^{e'}R(t,\Gamma)=\sum_i\hat h_i\Phi_i$ in $K[t,y]$, with
$c_2\ne0$. Hence $R(t,\Gamma(t,y))=0$ at every real nonsingular common zero
with $c_2(t)\ne0$. At a nonsingular zero with $c_2(t)=0$, the implicit
function theorem gives a branch of nonsingular zeros over a neighborhood of
$t$, on which $c_2\ne0$ on a dense set. Continuity extends the identity to the
given zero.
\end{proof}

\begin{lemma}[Ordered limits]
\label{lem:qc-ordered-limits}
Let $R\in K[\varepsilon,\delta,w]$ be nonzero. Suppose that
$R(\varepsilon_\ell,\delta_{\ell,\nu},w_{\ell,\nu})=0$ for nonzero reals
$\varepsilon_\ell\to0$ and $\delta_{\ell,\nu}\to0$, where
$w_{\ell,\nu}\to w_\ell$ as $\nu\to\infty$ and $w_\ell\to w^*$ as
$\ell\to\infty$, all limits finite. Then there is a nonzero $P\in K[w]$ with
$\deg P\le\deg_wR$ and $P(w^*)=0$.
\end{lemma}

\begin{proof}
Write $R=\delta^aS(\varepsilon,w)+\delta^{a+1}R_1$ with $S\ne0$, where $a$
is the least exponent of $\delta$ with a nonzero coefficient. Dividing by
$\delta_{\ell,\nu}^a$ and letting $\nu\to\infty$ gives
$S(\varepsilon_\ell,w_\ell)=0$ for every $\ell$; this remains true if
$S(\varepsilon_\ell,\cdot)$ is the zero polynomial. Write
$S=\varepsilon^bP(w)+\varepsilon^{b+1}S_1$ with $P\ne0$, divide by
$\varepsilon_\ell^b$, and let $\ell\to\infty$. Then $P(w^*)=0$. A nonzero
constant cannot vanish at $w^*$, so $P$ is not constant.
\end{proof}

We use the following classical root count.

\begin{theorem}[Multihomogeneous B\'ezout bound~\cite{MorganSommese1987,DedieuMalajovichShub2005}]
\label{thm:qc-bezout}
Let $\mathbb F$ be an algebraically closed field of characteristic zero and
$\Phi_1,\ldots,\Phi_N\in\mathbb F[Y,Z]$, where $Y$ and $Z$ are groups of $N_1$
and $N_2$ variables with $N_1+N_2=N$. If $\deg_Y\Phi_i\le e_i$ and
$\deg_Z\Phi_i\le e'_i$, then the number of isolated common zeros of $\Phi$ in
$\mathbb F^N$, and in particular the number of nonsingular ones, is at most the
coefficient of $U^{N_1}T^{N_2}$ in $\prod_{i=1}^N(e_iU+e'_iT)$.
\end{theorem}

The isolated-root upper bound is stated in
\cite[Section~5, Theorem~5.1]{DedieuMalajovichShub2005}.
Additional positive-dimensional components of the zero set are allowed in
this statement. For an algebraically closed field of characteristic zero and
countable transcendence degree, such as an algebraic closure of $K(t)$, the
statement follows from the case $\mathbb F=\C$ by embedding the field in $\C$.

\begin{theorem}[Joint degree]
\label{thm:qc-degree}
In the setting of this subsection,
\[
[K(p_1,\ldots,p_n):K]\le\mathcal B(n,k)=\max_{0\le s\le\min(k,n)}2^s\binom ns,
\]
and $\theta\in K(p_1,\ldots,p_n)$. For $K=\Q$, the value and the canonical
optimizer lie in one number field of degree at most $\mathcal B(n,k)\le
\min\{3^n,(2n)^k\}$.
\end{theorem}

\begin{proof}
Fix the nested sequences and the support $J$ of Lemma~\ref{lem:qc-kkt}. In
the system~\eqref{eq:qc-kkt}, each of the $d$ stationarity equations has degree
at most one in $u$ and at most one in $\lambda$, and each of the $s$ active
equations has degree at most two in $u$ and zero in $\lambda$. By
Theorem~\ref{thm:qc-bezout}, the number of nonsingular zeros over an algebraic
closure of $K(\varepsilon,\delta)$ is at most
\[
[U^dT^s](U+T)^d(2U)^s=2^s\binom ds .
\]
For $c\in K^n$, apply Lemma~\ref{lem:qc-parameter} with parameters
$t=(\varepsilon,\delta)$, unknowns $y=(u,\lambda)$, and output
$\Gamma=c^{\mathsf T}(x_0+Vu)$. Lemmas~\ref{lem:qc-ordered}
and~\ref{lem:qc-kkt} provide the nonsingular real zeros and the ordered
limits, and Lemma~\ref{lem:qc-ordered-limits} gives a nonzero
$P_c\in K[w]$ of degree at most $2^s\binom ds$ with $P_c(c^{\mathsf T}p)=0$.
Every coordinate of $p$ is therefore algebraic over $K$. The extension
$K(p)/K$ is finite and separable, so by the primitive element theorem some
$c\in K^n$ gives $K(p)=K(c^{\mathsf T}p)$. Hence
$[K(p):K]\le2^s\binom ds\le\mathcal B(n,k)$, using $s\le\min(k,d)$ and
$d\le n$. The value $\theta=f(p)$ lies in $K(p)$ because $f$ has coefficients
in $K$. Finally $2^s\binom ns\le(2n)^s$, and the sum over all $s$ is $3^n$.
\end{proof}

The same support $J$ serves all linear forms $c^{\mathsf T}p$. This is why
the bound concerns the joint field and not only individual coordinates;
separate coordinate bounds would not bound the degree of their compositum.
The maximum in $\mathcal B(n,k)$ is needed: for $n=k=3$ the terms are
$1,6,12,8$. The count $2^s\binom ds$ is the classical generic degree of
quadratically constrained critical points~\cite{NieRanestad2009}; what the
proof adds is its validity at the canonical optimizer of every convex
instance, through the span compression and the ordered limits, without
genericity or a constraint qualification.

\paragraph{The height bound.}
For a polynomial $\Psi$ with coefficients in $K$ and a place $v$ of $K$, let
$\norm\Psi_v$ be the sum of the $v$-absolute values of its coefficients if
$v$ is archimedean and their maximum otherwise. Both norms are
submultiplicative, and the second also satisfies the ultrametric inequality.

\begin{lemma}[Ordered elimination with local norms]
\label{lem:qc-elimination}
Let $G_1,\ldots,G_s,A,\Gamma\in K[\varepsilon,\delta,\lambda_1,\ldots,\lambda_s]$
have total degree at most $a\ge1$ in $\lambda$. For each place $v$ let
$C_v\ge1$ bound $\norm{G_i}_v$, $\norm A_v$ and $\norm\Gamma_v$, with $C_v=1$
for all but finitely many $v$, and put
\[
E=\frac1D\sum_vn_v\log C_v,\qquad \mathfrak z=(a+1)^s,\qquad T=a(s+1).
\]
Suppose that for nonzero reals $\varepsilon_\ell\to0$ and
$\delta_{\ell,\nu}\to0$ there are real roots $\lambda_{\ell,\nu}$ of
$G(\varepsilon_\ell,\delta_{\ell,\nu},\cdot)=0$ with nonsingular Jacobian in
$\lambda$ and $A\ne0$, such that $w_{\ell,\nu}=\Gamma/A$ at these roots
converges to $w_\ell$ as $\nu\to\infty$, and $w_\ell\to w^*$, all limits
finite. Then there is a nonzero $P\in K[w]$ with $\deg P\le\mathfrak z$, $P(w^*)=0$,
and
\[
\height_{\rm aff}(\text{coefficients of }P)\le W_0:=\mathfrak z(T+1)E+\log\mathfrak z!+\mathfrak z\log3 .
\]
\end{lemma}

\begin{proof}
Introduce a deformation parameter $\beta$ and the polynomials
$\Phi_i=G_i+\beta\lambda_i^{a+1}$. Over $K(\varepsilon,\delta,\beta)$, for a
monomial order that compares total degree in $\lambda$ first, the leading
monomials $\lambda_i^{a+1}$ are pairwise coprime, so the $\Phi_i$ form a
Gr\"obner basis and the monomials $\lambda^\alpha$ with all $\alpha_i\le a$
form a basis $\mathcal M$ of the quotient, with $\abs{\mathcal M}=\mathfrak z$. The same
holds after specializing $\varepsilon,\delta$ to real numbers and $\beta$ to a
nonzero real number, because the leading coefficients are $\beta$.

Let $M_A$ and $M_\Gamma$ be the matrices of multiplication by $A$ and
$\Gamma$ on this basis. A basis monomial times $A$ or $\Gamma$ has degree at
most $T$ in $\lambda$, and each replacement of $\lambda_i^{a+1}$ by
$-G_i/\beta$ lowers the degree by at least one. Hence at most $T$ replacements
occur along any branch of the reduction, the matrices $\beta^TM_A$ and
$\beta^TM_\Gamma$ have polynomial entries in $(\varepsilon,\delta,\beta)$, and
each entry has $v$-norm at most $C_v^{T+1}$: at an archimedean place each
replacement multiplies the coefficient sum by at most $C_v$, and at the other
places the maximum is multiplicative. Form
\[
\mathcal H=\det\bigl(w\beta^TM_A-\beta^TM_\Gamma-\zeta\beta^TI_{\mathfrak z}\bigr).
\]
Its coefficient of $\zeta^{\mathfrak z}$ is $(-\beta^T)^{\mathfrak z}\ne0$, its
degree in $w$ is at most $\mathfrak z$, and expanding the determinant gives
$\norm{\mathcal H}_v\le\mathfrak z!\,3^{\mathfrak z}C_v^{\mathfrak z(T+1)}$ at
archimedean places and $\norm{\mathcal H}_v\le C_v^{\mathfrak z(T+1)}$ at the
others.

Let $\mathcal H_\kappa$ be the coefficient of the least power $\zeta^\kappa$
occurring in $\mathcal H$. Fix one selected pair
$(\varepsilon_\ell,\delta_{\ell,\nu})$ and its root $\lambda_{\ell,\nu}$. The
real implicit function theorem continues the root to a real root
$\lambda(\beta)$ of $\Phi$ for small $\beta$, with $A\ne0$. For small
$\beta\ne0$, evaluation at $\lambda(\beta)$ is a ring homomorphism on the
quotient, so the vector of basis monomials evaluated at $\lambda(\beta)$ is a
nonzero common left eigenvector of the specialized multiplication matrices,
with eigenvalues $A(\lambda(\beta))$ and $\Gamma(\lambda(\beta))$. Hence the
specialized determinant equals $\beta^T(wA-\Gamma-\zeta)\,\Psi(w,\zeta)$ for
a polynomial $\Psi$. It is divisible by $\zeta^\kappa$, and $wA-\Gamma-\zeta$
is coprime to $\zeta$ because $A\ne0$, so $\zeta^\kappa$ divides $\Psi$.
Taking the coefficient of $\zeta^\kappa$ shows
$\mathcal H_\kappa(\varepsilon_\ell,\delta_{\ell,\nu},\beta,\Gamma/A)=0$ along
the branch, even if the specialized coefficient vanishes identically. This
step removes factors coming from roots at which $A=\Gamma=0$.

Let $\mathcal R(\varepsilon,\delta,w)$ be the coefficient of the least power
of $\beta$ in $\mathcal H_\kappa$. Dividing by that power and letting
$\beta\to0$ along the branch gives
$\mathcal R(\varepsilon_\ell,\delta_{\ell,\nu},w_{\ell,\nu})=0$ for all
selected pairs. Lemma~\ref{lem:qc-ordered-limits} now gives $P$ with
$P(w^*)=0$. Each extraction selects a subvector of coefficients, so no local
norm increases, and
\[
\height_{\rm aff}(\text{coefficients of }P)\le
\frac1D\sum_vn_v\log\max\bigl(1,\norm{\mathcal H}_v\bigr)\le W_0,
\]
because the archimedean weights sum to $D$.
\end{proof}

Deforming a polynomial system to obtain a finite quotient, and reading the
values of a function from multiplication matrices, are classical elimination
devices~\cite{Canny1990,GrigorievPasechnik2005}. The lemma records the
explicit bound, place by place, that the following theorem needs.

\begin{theorem}[Height and separation]
\label{thm:qc-height}
There is an absolute constant $c$ such that, in the setting of this
subsection, with $W=(B+1)S^{c(k+1)}$:
\begin{enumerate}[label=(\roman*),leftmargin=*]
\item $\height(\theta)\le W$ and $\height(p_j)\le W$ for every $j$; more
generally $\height(c^{\mathsf T}p)\le W$ for every $c\in K^n$ whose entries
have height at most $B$;
\item if $\theta\ne0$, then $\abs{\log\abs\theta}\le DW$; the same holds for
every nonzero $p_j$;
\item for $K=\Q$, the primitive minimal polynomials of $\theta$, of each
$p_j$, and of each $c^{\mathsf T}p$ as in (i) have degree at most
$\mathcal B(n,k)$ and coefficient bit length at most $L^{c(k+1)}$, and every
nonzero one of these numbers lies in
$[2^{-L^{c(k+1)}},2^{L^{c(k+1)}}]$ in absolute value.
\end{enumerate}
\end{theorem}

\begin{proof}
Assume $d\ge1$. Fix $J$ and the sequences of Lemma~\ref{lem:qc-kkt}, and write
$M(\varepsilon,\lambda)=\tilde Q_0+2\varepsilon V^{\mathsf T}V+
\sum_{j\in J}\lambda_j\tilde Q_j$ and
$\mu(\varepsilon,\lambda)=\tilde a_0+2\varepsilon V^{\mathsf T}x_0+
\sum_{j\in J}\lambda_j\tilde a_j$, so that the stationarity equations
in~\eqref{eq:qc-kkt} read $Mu+\mu=0$. Let $\Delta=\det M$ and
$\rho=-\operatorname{adj}(M)\mu$. At the selected roots $M\succ0$ and
$u=\rho/\Delta$. Put
\[
G_j=\tfrac12\rho^{\mathsf T}\tilde Q_j\rho+\Delta\tilde a_j^{\mathsf T}\rho
+(\tilde c_j-\delta)\Delta^2\quad(j\in J),\qquad A=\Delta^2,
\]
and, for the outputs, $\Gamma_f=\tfrac12\rho^{\mathsf T}\tilde Q_0\rho+
\Delta\tilde a_0^{\mathsf T}\rho+\tilde c_0\Delta^2$ and
$\Gamma_c=\Delta(c^{\mathsf T}x_0\Delta+c^{\mathsf T}V\rho)$ for $c\in K^n$.
Then $G_j=\Delta^2(r_j(u)-\delta)$, $\Gamma_f/A=\tilde f(u)$ and
$\Gamma_c/A=c^{\mathsf T}x$. Along the root branch
$u(\lambda)=-M^{-1}\mu$, one has $\partial u/\partial\lambda_i=
-M^{-1}\nabla r_i(u)$, and because $r_j(u)=\delta$ at the root,
\[
\frac{\partial G_j}{\partial\lambda_i}=
-\Delta^2\,\nabla r_j(u)^{\mathsf T}M^{-1}\nabla r_i(u),
\]
which is negative definite as a matrix in $(j,i)$. So the selected multiplier
vectors are nonsingular roots of $G$ with $A\ne0$. The ordered limits of
Lemma~\ref{lem:qc-ordered} give $\Gamma_f/A\to\theta$ and
$\Gamma_c/A\to c^{\mathsf T}p$.
All these polynomials have degree at most $a=2d$ in $\lambda$.

Let $\mathcal S$ be the list of scalars consisting of the entries of
$\tilde Q_i,\tilde a_i,\tilde c_i$ for $i\in\{0\}\cup J$, of
$2V^{\mathsf T}V$, $2V^{\mathsf T}x_0$, $x_0$, $V$, $c^{\mathsf T}x_0$ and
$c^{\mathsf T}V$, and the constants $1,2,\tfrac12$. For a place $v$ let $U_v=\max(1,\max_{\sigma\in\mathcal S}
\abs\sigma_v)$. At an archimedean place each entry of $M$ and $\mu$ has norm
at most $(s+2)U_v$, so $\Delta$ and every entry of $\rho$ have norm at most
$\Pi_v=d!((s+2)U_v)^d$, and every $G_j$, $A$, $\Gamma_f$, $\Gamma_c$ has norm
at most $(d+1)^2U_v^2\Pi_v^2$. At a nonarchimedean place the same polynomials
have norm at most $U_v^{2d+2}$. Hence, with
$H_{\mathcal S}=\sum_{\sigma\in\mathcal S}\height(\sigma)$,
\[
E\le(2d+2)H_{\mathcal S}+2\log((d+1)!)+2d\log(s+2).
\]
By Lemma~\ref{lem:qc-restriction}(d), $\abs{\mathcal S}\le S^{c_0+3}$ and
$H_{\mathcal S}\le S^{2c_0+3}(B+1)$. Lemma~\ref{lem:qc-elimination}, with
$\mathfrak z=(2d+1)^s$ and $T=2d(s+1)$, gives annihilators $P$ of $\theta$ and of each
$c^{\mathsf T}p$ with
\[
\height_{\rm aff}(\text{coefficients of }P)\le
(2d+1)^s\bigl[(2d(s+1)+1)E+s\log(2d+1)+\log3\bigr].
\]
Since $d\le S$ and $s\le k$, this is at most $(B+1)S^{c(k+1)}$ for an
absolute $c$.

Let $\vartheta$ be a root of $P=\sum_i\pi_iw^i$ of degree $e$. By Fujiwara's
bound at archimedean places and the ultrametric bound at the others,
$\log\max(1,\abs\vartheta_u)\le\log2\,[u\mid\infty]+
\log\max_i\abs{\pi_i}_u-\log\abs{\pi_e}_u$ at every place $u$ of a field
containing $\vartheta$. Summing with the weights and using the product formula
for $\pi_e$ gives $\height(\vartheta)\le\height_{\rm proj}(P)+\log2$, which
proves (i) after enlarging $c$. For (ii), remove a power of $w$ from $P$ so
that its constant coefficient $\pi'_0$ is nonzero. Each ratio $\pi'_i/\pi'_0$ lies
in $K$ and has height at most $\height_{\rm proj}(P)$ by
Lemma~\ref{lem:qc-heights}(a), hence absolute value at most
$\exp(D\height_{\rm proj}(P))$ by Lemma~\ref{lem:qc-heights}(c). Fujiwara's
bound applied to $P$ and to its reversal gives
$\abs{\log\abs\vartheta}\le D\height_{\rm proj}(P)+\log2$. For $K=\Q$,
clearing denominators turns $P$ into a primitive integer polynomial whose
largest coefficient has logarithm $\height_{\rm proj}(P)$. Then
Theorem~\ref{thm:qc-degree} and~\eqref{eq:qc-minpoly-height} give (iii), with
$S,B\le L$ and $\mathcal B(n,k)\le(2L)^k$.
\end{proof}

The degree factor $D$ enters (ii) linearly. The elimination is carried out
over $K$ with local norms, so the field is never encoded by a quantified
primitive element, and neither $D$ nor $B$ is raised to a power depending on
$k$. Convexity is used only at the fixed real embedding; the conjugate
problems need not be convex. Theorem~\ref{thm:qc-degree} applies over $K$ as
well and bounds $[K(\theta):K]$ by $\mathcal B(n,k)$, so
$[\Q(\theta):\Q]\le D\,\mathcal B(n,k)$.

\begin{corollary}[Radius and value gap]
\label{cor:qc-radius}
There is an absolute constant $c$ with the following properties for every
rational native PSD system with Hessian span $k$ and input length $L$.
\begin{enumerate}[label=(\alph*),leftmargin=*]
\item If $F\ne\varnothing$, its minimum-norm point lies in $[-R_F,R_F]^n$,
where $R_F=2^{\lceil L^{c(k+1)}\rceil}$.
\item Let a rational box $[-R,R]^n$, a rational affine function $\phi$, and
rational affine rows be added, and let $\phi^*$ be the minimum of $\phi$ over
the resulting compact set, assumed nonempty. If $\phi^*\ne0$, then
$\abs{\phi^*}\ge2^{-(L+\log_2R)L^{c(k+1)}}$, where $L$ now counts the added
rows as well.
\end{enumerate}
\end{corollary}

\begin{proof}
For (a), apply Theorem~\ref{thm:qc-height}(ii) with $f=0$. For (b), the
minimum is attained by compactness, the added rows are affine, and
Theorem~\ref{thm:qc-height}(ii) with $D=1$ applies with $B\le L+\log_2R+2$.
\end{proof}

The bound in (b) is linear in $\log R$. Boxes of radius $R_F$ therefore add
only a factor $L^{O(k+1)}$ to the exponent, and the gap remains
$2^{-L^{O(k+1)}}$.

\begin{proposition}[Exact feasibility at fixed span]
\label{prop:qc-decision}
For each fixed $k$, a deterministic algorithm decides in time polynomial in
$L$ whether a rational native PSD system with Hessian span at most $k$ is
feasible. Its running time is $L^{O(k+1)}$.
\end{proposition}

\begin{proof}
Compute $R=R_F$ from Corollary~\ref{cor:qc-radius}(a). Let
$\nu(x)=\max\{0,g_1(x),\ldots,g_m(x),(Ax-b)_1,\ldots\}$, a convex function
with exact rational values and subgradients at rational points, and let
$\alpha=\min\{\nu(x):\norm x_\infty\le R\}$. If $F\ne\varnothing$, its
minimum-norm point gives $\alpha=0$; if $\alpha=0$, a minimizer lies in $F$.
The number $\alpha$ is the optimal value of the epigraph problem
$\min\{\tau:g_i(x)\le\tau,\ (Ax-b)_j\le\tau,\ 0\le\tau,\ \norm x_\infty\le R\}$,
whose constraint Hessians $\diag(Q_i,0)$ span $k$ dimensions; adding the row
$\tau\le\bar\nu$, where $\bar\nu$ is an explicit rational upper bound for
$\nu$ on the box, makes its feasible set compact without changing $\alpha$. By
Corollary~\ref{cor:qc-radius}(b), $\alpha=0$ or $\alpha\ge\eta:=
2^{-L^{c'(k+1)}}$ for an absolute $c'$. Termwise estimates give rational
bounds of polynomial length for $\abs\nu$ and for the norms of the gradients
of the rows on the box, and a subgradient of $\nu$ at a rational point is the
gradient of a row attaining the maximum, or zero. Hence
Lemma~\ref{lem:convex-value}(a), with accuracy $\eta/3$, returns a rational
point $\hat x$ of the box with $\nu(\hat x)\le\alpha+\eta/3$, in time
polynomial in $n$, $\log R$, $\log(1/\eta)$ and the input length. The system
is feasible exactly when $\nu(\hat x)<\eta/2$.
\end{proof}

This is the exact decision procedure of the Hessian-span program. Its
exponent grows with $k$; no running time of the form $\phi(k)L^{O(1)}$ is
claimed.

\paragraph{Sharpness.}
The degree bound is attained, even under strong regularity. The proof
specializes a generic count inside a nonempty open set of strictly convex
instances, using two classical inputs. Nie and Ranestad prove that for
generic coefficients the equality-constrained critical point system
\begin{equation}
\label{eq:qc-generic-kkt}
g_i(x)=0\ (1\le i\le s),\qquad
\nabla f(x)+\sum_{i=1}^s\lambda_i\nabla g_i(x)=0
\end{equation}
of quadratics in $n$ variables has exactly $2^s\binom ns$ complex
solutions~\cite{NieRanestad2009}. Hilbert's irreducibility theorem, in the
form that the integral specializations at which an irreducible polynomial over
$\Q(u_1,\ldots,u_r)$ becomes reducible or drops degree form a set of density
zero in boxes of $\Z^r$, supplies the rational
specialization~\cite{Serre2008}.

\begin{theorem}[Sharpness of the degree bound]
\label{thm:qc-sharp}
For every $n\ge1$ and $0\le k\le n(n+1)/2$ there is a rational instance with
Hessian span exactly $k$, positive definite objective Hessian, a strictly
feasible compact feasible set, and a unique optimizer $p$, such that
$[\Q(p):\Q]=[\Q(\theta):\Q]=\mathcal B(n,k)$. If $k>0$, all constraint
Hessians can be taken positive definite.
\end{theorem}

\begin{proof}
For $k=0$ minimize $\norm x^2$ over a box. Let $k\ge1$ and choose
$1\le s\le\min(k,n)$ attaining $\mathcal B(n,k)$; the term for $s=1$ is $2n\ge1$,
so such $s$ exists. Let $t$ be the vector of all coefficients of
$f,g_1,\ldots,g_s$, treated as indeterminates, $\mathbb K=\Q(t)$, and
$\mathcal D=2^s\binom ns$. The incidence variety of~\eqref{eq:qc-generic-kkt}
is the graph of a polynomial map: $x$, $\lambda$, all Hessians,
$a_1,\ldots,a_s$ and $c_0$ are free, and $c_i=-\tfrac12x^{\mathsf T}Q_ix-
a_i^{\mathsf T}x$ and $a_0=-Q_0x-\sum_i\lambda_i(Q_ix+a_i)$ are determined.
It is therefore an irreducible variety, of the same dimension as the
parameter space. Its projection to the parameter space has generic fibers of
cardinality $\mathcal D$, so its function field $\mathcal L=\mathbb K(x,\lambda)$
is a separable extension of degree $\mathcal D$ of $\mathbb K$. Cramer's rule
on a nonzero minor of the gradient matrix gives $\lambda\in\mathbb K(x)$, so
$\mathcal L=\mathbb K(x)$.

Let $\beta=f(x)\in\mathcal L$. Each derivation $\partial_j=\partial/\partial
a_{0,j}$ of $\mathbb K$ extends uniquely to $\mathcal L$. Differentiating
$g_i(x)=0$ gives $\nabla g_i(x)^{\mathsf T}\partial_jx=0$, and therefore
\[
\partial_j\beta=x_j+\nabla f(x)^{\mathsf T}\partial_jx
=x_j-\sum_i\lambda_i\nabla g_i(x)^{\mathsf T}\partial_jx=x_j .
\]
A derivation of $\mathbb K$ maps $\mathbb K(\beta)$ into itself: if
$\Pi$ is the minimal polynomial of $\beta$ over $\mathbb K$, then
$\partial_j\beta=-(\partial_j\Pi)(\beta)/\Pi'(\beta)$. Hence
$x_j\in\mathbb K(\beta)$ and $\mathbb K(\beta)=\mathcal L$. Write
$x_j=\varphi_j(\beta)$ and $\lambda_i=\psi_i(\beta)$ with
$\varphi_j,\psi_i\in\mathbb K[Y]$, and let $\mathfrak c\in\Q[t]$ be a nonzero
common denominator of these polynomials and of $\Pi$. The identities
$\Pi(\beta)=0$, $x_j=\varphi_j(\beta)$ and $\lambda_i=\psi_i(\beta)$ hold on
the incidence variety wherever $\mathfrak c\ne0$.

For the seed $Q_i=I+e_ie_i^{\mathsf T}$, $a_i=e_i$, $c_i=0$ ($1\le i\le s$),
$Q_0=I$, $a_0=-\sum_ie_i$, $c_0=0$, the point $x=0$ with $\lambda=\mathbf1$
solves~\eqref{eq:qc-generic-kkt}, all Hessians are positive definite,
$Q_1,\ldots,Q_s$ are independent, and the active gradients $e_1,\ldots,e_s$ are
independent. The bordered Jacobian is nonsingular, so the implicit function
theorem gives a real solution branch on a neighborhood of the seed. Shrinking
the neighborhood, we obtain a nonempty open set $\mathcal O$ of real
coefficient vectors on which the Hessians are positive definite and
independent, the multipliers are positive, a fixed point is strictly feasible
for $g_i\le0$, and the branch point is, by convexity, the unique optimizer of
$\min\{f:g_i\le0\}$. Choose a rational $t_0\in\mathcal O$ and put
$t=t_0+(1/u_1,\ldots,1/u_r)$. This birational substitution keeps $\Pi$
irreducible over $\Q(u)$. By Hilbert's theorem in the form above, there are
integral $u$ with all $\abs{u_j}$ so large that $t\in\mathcal O$, at which
$\Pi$ stays irreducible of degree $\mathcal D$ and the leading coefficient and
$\mathfrak c$ do not vanish: each excluded condition defines a set of density
zero.

At such a rational $t$, the unique optimizer $p$ is a point of the incidence
variety, so $\beta^*=f(p)$ is a root of the irreducible specialized $\Pi$ and
$[\Q(\beta^*):\Q]=\mathcal D$. The specialized identities give
$p\in\Q(\beta^*)^n$, and $\beta^*\in\Q(p)$, so $\Q(p)=\Q(\beta^*)$ has degree
$\mathcal D=\mathcal B(n,k)$. If $s<k$, extend $Q_1,\ldots,Q_s$ to $k$
independent rational positive definite matrices, which is possible because
the positive definite cone is open and spans the symmetric matrices, and add
the rows $\tfrac12x^{\mathsf T}Rx-C\le0$ with rational $C$ so large that they
are strictly slack at $p$ and at a strictly feasible point. The optimizer and
value do not change, and the Hessian span becomes exactly $k$.
\end{proof}

The theorem shows that degeneracy, lack of strict feasibility, and
nonconvexity are not needed to reach the worst degree. It is an existence
statement: it gives no useful bound on the coefficient sizes of the
instances. Section~\ref{app:qc-short} constructs explicit short instances with
degree $n^{\Omega(k)}$.

\subsection{Rational points at span two and irrational singletons at span three}
\label{app:qc-rational}

A rational feasible point is the simplest exact output, and it is what an
exact solver would like to return. For one convex quadratic row together with
rational affine rows, a nonempty system has a rational point of polynomial
length~\cite{Vavasis1990,DelPiaDeyMolinaro2017}. This subsection shows that
the same is true whenever the native Hessians span at most two dimensions, and
that three dimensions already permit a feasible set consisting of one
irrational point. The threshold is therefore exactly between two and three.
The difficult case is a tangency between two convex quadratics, and the key
fact is that a tangency multiplier is rational.

\begin{lemma}[Rational tangency multiplier]
\label{lem:qc-tangency}
Let $g_1,g_2$ be convex quadratics on $\R^d$ with rational coefficients,
$g_i(x)=\tfrac12x^{\mathsf T}Q_ix+a_i^{\mathsf T}x+c_i$. Suppose that for
some $t_0>0$ and $z\in\R^d$,
\[
g_1+t_0g_2\ge0\ \text{on }\R^d,\qquad g_1(z)=g_2(z)=0 .
\]
Then there is a rational $t>0$, of bit length polynomial in the input length,
with $g_1+tg_2\ge0$ on $\R^d$ and $(g_1+tg_2)(z)=0$. It is computed in
polynomial time by the following rule. Let $W=\ker Q_1\cap\ker Q_2$. If some
rational $w\in W$ has $(a_1^{\mathsf T}w,a_2^{\mathsf T}w)\ne(0,0)$, take
$t=-a_1^{\mathsf T}w/a_2^{\mathsf T}w$. Otherwise restrict to a rational
complement of $W$ and form
\[
\varphi(\tau)=c_1+\tau c_2-\tfrac12(a_1+\tau a_2)^{\mathsf T}
(Q_1+\tau Q_2)^{-1}(a_1+\tau a_2)\in\Q(\tau);
\]
take $t=1$ if $\varphi=0$, and otherwise let $t$ be the root of
$\gcd(\mathcal N,\mathcal N')$, where $\mathcal N$ is the numerator of
$\varphi$ in lowest terms.
\end{lemma}

\begin{proof}
The aggregate $g=g_1+t_0g_2$ is nonnegative and vanishes at $z$, so
$\nabla g(z)=0$. For $w\in W$, $\nabla g(z)^{\mathsf T}w=(a_1+t_0a_2)^{\mathsf
T}w$, hence $(a_1+t_0a_2)^{\mathsf T}w=0$. If $(a_1^{\mathsf T}w,a_2^{\mathsf
T}w)\ne(0,0)$ for some rational $w\in W$, then $a_2^{\mathsf T}w\ne0$ and
$t_0=-a_1^{\mathsf T}w/a_2^{\mathsf T}w$ is rational; take $t=t_0$.

Otherwise both linear forms vanish on $W$, and $g_1,g_2$ are constant along
$W$. Pass to the rational complement $W^\perp$, on which $Q_1+Q_2\succ0$ and
hence $Q_1+\tau Q_2\succ0$ for $\tau>0$. If $W^\perp=0$, both quadratics are
constants that vanish at $z$, and $t=1$ works. Otherwise $\varphi(\tau)$ is
the minimum of $g_1+\tau g_2$ for $\tau>0$, and $\varphi(t_0)=0$. The minimizer
at $t_0$ is the image $\bar z$ of $z$, and the derivative of the minimum value
with respect to $\tau$ is $g_2$ at the minimizer, so $\varphi'(t_0)=g_2(\bar z)
=0$. If $\varphi=0$, then $g_1+g_2\ge0$ with minimum $\varphi(1)=0$, and
$g_1(z)+g_2(z)=0$, so $t=1$ works.

Suppose $\varphi\ne0$. A real congruence brings $Q_1+Q_2$ to the identity and
$Q_2$ to $\diag(d_1,\ldots,d_{d'})$ with $0\le d_j\le1$, hence $Q_1$ to
$\diag(1-d_j)$. Dividing each scalar term gives the partial-fraction form
\begin{equation}
\label{eq:qc-partial-fractions}
\varphi(\tau)=\kappa_0+\kappa_1\tau-\kappa_2\tau^2-\sum_{j=1}^r\frac{\rho_j}{\tau-\sigma_j},
\qquad \kappa_2\ge0,\quad\rho_j>0,\quad\sigma_1<\cdots<\sigma_r\le0 .
\end{equation}
Terms with $d_j=0$ contribute to the quadratic polynomial with nonpositive
leading coefficient; a term with $d_j>0$ has its pole at
$\sigma=-(1-d_j)/d_j\le0$ with a nonpositive residue. Equal poles are
combined and zero residues removed, so the $\sigma_j$ are exactly the poles of
$\varphi$, all simple. The diagonalization is used only to obtain these signs;
the numbers in~\eqref{eq:qc-partial-fractions} need not be rational. Write
$\varphi=\mathcal N/\mathcal Q$ in lowest terms over $\Q$. Then
$\deg\mathcal N\le r+1$ if $\kappa_2=0$, and $\deg\mathcal N=r+2$ if
$\kappa_2>0$. On each interval $(\sigma_j,\sigma_{j+1})$, $\varphi$ tends to
$-\infty$ at the left end and to $+\infty$ at the right end, so $\mathcal N$ has
a root of odd multiplicity there. If $\kappa_2>0$, there is one more such root
in $(-\infty,\sigma_1)$. All these roots are nonpositive and distinct from
$t_0>0$, while $t_0$ is a root of multiplicity at least two because
$\varphi(t_0)=\varphi'(t_0)=0$ and $\mathcal Q(t_0)\ne0$. Counting with
multiplicity, these roots exhaust $\deg\mathcal N$ when $r\ge1$. Hence every
other root of $\mathcal N$ is simple and $t_0$ is a root of multiplicity
exactly two. When $r=0$, $\varphi$ is a polynomial of degree at most two; it
cannot be affine with a double root, so it equals $-\kappa_2(\tau-t_0)^2$. In
all cases $\gcd(\mathcal N,\mathcal N')$ has degree one and root $t_0$, so
$t_0\in\Q$.

A numerator of $\varphi$ before cancellation is
\[
2(c_1+\tau c_2)\det(Q_1+\tau Q_2)-(a_1+\tau a_2)^{\mathsf T}
\operatorname{adj}(Q_1+\tau Q_2)(a_1+\tau a_2),
\]
which has degree at most $d+1$ and coefficients of polynomial bit length.
The rational root theorem bounds the numerator and denominator of $t_0$. The
determinant, the adjugate, the cancellation and the gcd are polynomial-time
rational computations.
\end{proof}

Cancellation must precede the gcd. In $\R^3$, the balls
$g_1=\norm x^2-1$ and $g_2=\norm{x-2e_1}^2-1$ touch at $e_1$. Here
$\varphi(\tau)=-(\tau-1)^2/(\tau+1)$, whereas the uncancelled numerator above
is $-16(\tau+1)^2(\tau-1)^2$, with a spurious repeated root $-1$. Repeated
roots of the characteristic determinant also occur in classical contact tests
for ellipsoids~\cite{JiaChoiMourrainWang2011}; only the reduced numerator has
the uniqueness property used in the proof.

\begin{proposition}[Two quadratic rows]
\label{prop:qc-two-rows}
Let $P$ be a rational polyhedron and $g_1,g_2$ rational convex quadratics. If
$F=P\cap\{g_1\le0,g_2\le0\}$ is nonempty, then $F$ contains a rational point.
\end{proposition}

\begin{proof}
Let $I_t$ be the set of affine rows of $P$ that vanish identically on $F$, and
$\mathcal Y$ the rational affine space defined by them. For each other row choose
a point of $F$ at which it is strict, and let $p_0$ be the average of these
points, or any point of $F$ if there are none. Then $p_0\in F$ and every row
outside $I_t$ is strict at $p_0$.

If some $y\in\mathcal Y$ has $g_1(y)<0$ and $g_2(y)<0$, the points
$p_0+s(y-p_0)$ for small $s>0$ lie in $\mathcal Y$, satisfy every row outside
$I_t$ strictly, and make both quadratics negative. A neighborhood of such a
point in $\mathcal Y$ lies in $F$, and rational points are dense in the rational
affine space $\mathcal Y$.

Otherwise the convex open set
$\{(u_1,u_2):\exists x\in\mathcal Y,\ g_1(x)<u_1,\ g_2(x)<u_2\}$ misses the
closed negative quadrant. A separating line gives weights
$\lambda_1,\lambda_2\ge0$, not both zero, with $\lambda_1g_1+\lambda_2g_2\ge0$
on $\mathcal Y$; the weights are nonnegative because the open set is closed under
increasing $u$, and the separating constant is zero because $p_0$ is feasible.
Work in a rational affine chart of $\mathcal Y$.

If only $\lambda_1>0$, then $g_1\ge0$ on $\mathcal Y$ and $g_1(p_0)\le0$, so
every point of $F$ minimizes $g_1$ on $\mathcal Y$. These minimizers form the
rational affine space where the gradient of $g_1$ on $\mathcal Y$ vanishes.
Restricting to it leaves one quadratic row and affine rows. Repeating the
argument with this single row gives either a strict point, hence a rational
point by density, or a second restriction to the zero set of that row, after
which only affine rows remain and $F$ is a nonempty rational polyhedron. The
case where only $\lambda_2>0$ is symmetric.

If both weights are positive, then $g_1(p_0)=g_2(p_0)=0$, and
Lemma~\ref{lem:qc-tangency} with $t_0=\lambda_2/\lambda_1$ gives a rational
$t>0$ for which $g=g_1+tg_2$ is nonnegative on $\mathcal Y$ and vanishes at
$p_0$. Every point of $F$ has $g\le0$ and therefore lies in the zero set $Z$
of $g$ on $\mathcal Y$, which is the rational affine space where its gradient
vanishes. The direction space of $Z$ is the kernel of the restricted Hessian
of $g$, which equals the intersection of the kernels of the restricted $Q_1$
and $Q_2$ because both are positive semidefinite. Hence $g_1$ and $g_2$ are
affine on $Z$, and $F=Z\cap P\cap\{g_1\le0,g_2\le0\}$ is a nonempty rational
polyhedron.
\end{proof}

\begin{theorem}[Span two]
\label{thm:qc-span-two}
Let $F$ be a rational native PSD system with Hessian span $k\le2$.
\begin{enumerate}[label=(\alph*),leftmargin=*]
\item If $F\ne\varnothing$, then rational points are dense in $F$ and the
affine hull of $F$ is a rational affine space.
\item A deterministic polynomial-time algorithm either reports that
$F=\varnothing$ or returns a rational point of $F$. In particular, every
nonempty $F$ has a rational point of polynomial bit length.
\end{enumerate}
No boundedness, strict feasibility, positive definiteness, or bound on the
number of rows is assumed.
\end{theorem}

\begin{proof}
First reduce to at most two quadratic rows. If $k=2$, the cone generated by
$Q_1,\ldots,Q_m$ lies in a two-dimensional space and is pointed, because a
matrix and its negative cannot both be nonzero and positive semidefinite. Its
two extreme rays are generated by input matrices $B_1,B_2$, so
$Q_i=\alpha_iB_1+\beta_iB_2$ with rational $\alpha_i,\beta_i\ge0$. The system
is the projection of
\[
\tfrac12x^{\mathsf T}B_1x\le u_1,\quad\tfrac12x^{\mathsf T}B_2x\le u_2,\quad
\alpha_iu_1+\beta_iu_2+a_i^{\mathsf T}x+c_i\le0\ (1\le i\le m),\quad Ax\le b:
\]
take $u_j=\tfrac12x^{\mathsf T}B_jx$ in one direction and use
$\alpha_i,\beta_i\ge0$ in the other. Rational lifted points project to
rational points. The rays can be found by testing at most $m^2$ pairs. For
$k=1$ one generator suffices, and for $k=0$ the system is polyhedral.

(a) Proposition~\ref{prop:qc-two-rows} gives a rational point. Intersecting
$F$ with small rational boxes around any feasible point preserves the span
and nonemptiness, which proves density. Rational feasible points close to an
affine basis of the affine hull remain affinely independent, so the affine
hull is spanned by rational points.

(b) Work with the lifted system, which has span at most two and polynomial
length; call its feasible set $F$ again, its dimension $n$, and its
quadratic rows $g_1,g_2$. Every decision below concerns a rational native PSD
system of span at most two and polynomial length, obtained by adding rational
affine rows, so Proposition~\ref{prop:qc-decision} answers it in polynomial
time.

\emph{Step 1.} Decide whether $F=\varnothing$. If not, let $R=R_F$ from
Corollary~\ref{cor:qc-radius}(a), so $F$ meets $[-R,R]^n$, and let
$\mathrm{Box}=[-2R,2R]^n$.

\emph{Step 2: the affine hull of the minimal face.} For each affine row
$\ell_j\le0$, let $\eta_j$ be the gap of Corollary~\ref{cor:qc-radius}(b) for
minimizing $\ell_j$ over $F\cap\mathrm{Box}$, and decide whether
$F\cap\mathrm{Box}\cap\{\ell_j\le-\eta_j\}\ne\varnothing$. This happens exactly
when $\ell_j$ is not identically zero on $F$: if $\ell_j(y)<0$ for some
$y\in F$, the segment from a point of $F\cap[-R,R]^n$ toward $y$ contains
points of $F\cap\mathrm{Box}$ with $\ell_j<0$, so the minimum over the compact
set $F\cap\mathrm{Box}$ is negative and hence at most $-\eta_j$. Let $I_t$ be
the rows that are identically zero on $F$, $\mathcal Y$ the affine space where
they vanish, and $x=x_0+Uv$ a rational chart of $\mathcal Y$ whose free
coordinates are original coordinates. If $\mathcal Y$ is a point, return it.
Averaging the points found for the rows outside $I_t$ gives, in the proof
only, a point $\bar x\in F\cap\mathrm{Box}$ at which every row outside $I_t$
is strictly negative.

\emph{Step 3: a common margin.} Let $\eta$ be the gap of
Corollary~\ref{cor:qc-radius}(b) for maximizing $\sigma$ subject to
$g_i\le-\sigma$, $\ell_j\le-\sigma$ for $j\notin I_t$, $0\le\sigma\le1$, and
$v\in[-3R,3R]^{\dim\mathcal Y}$. Decide whether the set $C_\eta$ of chart points
in this box with $g_i\le-\eta$ and $\ell_j\le-\eta$ for $j\notin I_t$ is
nonempty. If some $y\in\mathcal Y$ has $g_1(y),g_2(y)<0$, the segment from
$\bar x$ toward $y$ gives a point of positive margin in the box, so
$C_\eta\ne\varnothing$; the converse is clear.

\emph{Step 4: positive margin.} If $C_\eta\ne\varnothing$, bisect the box one
coordinate at a time, keeping a closed half whose intersection with $C_\eta$
is nonempty; one decision per bisection suffices, since the other half is
nonempty when the tested half is empty. Stop when all side lengths are at
most $2\omega$, where $\omega=\min(1,\eta/(2G))$ and $G\ge1$ is a rational
Lipschitz bound in the maximum norm for the $g_i$ and $\ell_j$ on the box
enlarged by one. The center $\hat v$ is within $\omega$ of a point of
$C_\eta$, so every quadratic and every row outside $I_t$ is at most
$-\eta/2<0$ at $\hat v$, and $x_0+U\hat v$ is a rational point of $F$. The
number of bisections is polynomial.

\emph{Step 5: zero margin.} If $C_\eta=\varnothing$, no point of $\mathcal Y$
makes both quadratics negative. Compute the infimum of each $g_i$ on $\mathcal Y$
by rational linear algebra. If it is zero for some $i$, every feasible point
lies in the rational affine space where the gradient of $g_i$ on $\mathcal Y$
vanishes, and $g_i=0$ there; restrict to that space, delete $g_i$, and repeat
Steps 2--5 with one quadratic row. With one row and zero margin, the row is
nonnegative on the current affine space, so its infimum is zero and the next
restriction leaves a rational polyhedron. If both infima are negative, the
weights in the convex alternative of Proposition~\ref{prop:qc-two-rows} are
both positive, since a single positive weight would make the corresponding
quadratic nonnegative on $\mathcal Y$. Lemma~\ref{lem:qc-tangency} then computes
a rational $t>0$, and the rational affine space where the gradient of
$g_1+tg_2$ on $\mathcal Y$ vanishes contains $F$ and carries both quadratics as
affine functions.

In both branches a nonempty rational polyhedron remains, and rational linear
programming returns a rational point in polynomial
time~\cite{GroetschelLovaszSchrijver1988}. At most two curved restrictions
occur, so the polynomial bounds compose a bounded number of times. The output
of a polynomial-time algorithm has polynomial length.
\end{proof}

The face restriction in Step 2 is needed before a zero margin can be read as
a tangency. For $g_1=x^2-1$, $g_2=y^2-1$ and the row $x\ge1$, there is no
common positive margin, yet both quadratics are negative at the origin and no
nonzero nonnegative aggregate is nonnegative on the plane. Step 2 restricts to
$x=1$, where $g_1$ vanishes identically.

Rational points of convex semialgebraic sets can be computed in time
exponential in the dimension~\cite{SafeyElDinZhi2010}; for native PSD systems
of Hessian span at most two, the algorithm above is polynomial in the
dimension and the number of rows. The theorem gives a rational point, not the
canonical one. The minimum-norm
point of $\{x:(x-2)^2\le2\}$ is $2-\sqrt2$, although $1$ is feasible. Two
consequences are worth recording. First, if a rational convex quadratic
objective has a rational attained optimal value $\theta$, and the objective and
constraint Hessians together span at most two dimensions, appending
$f\le\theta$ yields a rational optimizer in polynomial time. Second, in a
mixed-integer system whose continuous Hessian span is at most two, every
nonempty continuous fiber over a fixed integer assignment has a rational
point; the theorem says nothing about finding the assignment.

\begin{proposition}[An irrational singleton at span three]
\label{prop:qc-span-three}
Let $r=\sqrt[3]2$ and
\[
g_1=x^2-y,\qquad g_2=y^2-2x,\qquad g_3=(x-y)^2-2x-y+4 .
\]
These rows have linearly independent rank-one positive semidefinite Hessians,
and $\{g_1\le0,g_2\le0,g_3\le0\}=\{(r,r^2)\}$. The rows
$R_i=3g_i+g_1+g_2+g_3$ have positive definite Hessians, Hessian span three,
and the same feasible set. No nonzero rational weights $w\ge0$ make
$\sum_iw_ig_i$ nonnegative on $\R^2$.
\end{proposition}

\begin{proof}
The Hessians are, up to the factor two,
$\left(\begin{smallmatrix}1&0\\0&0\end{smallmatrix}\right)$,
$\left(\begin{smallmatrix}0&0\\0&1\end{smallmatrix}\right)$ and
$\left(\begin{smallmatrix}1&-1\\-1&1\end{smallmatrix}\right)$, which are
independent. With $p=(r,r^2)$ and $r^3=2$, all three rows vanish at $p$; for
instance $g_3(p)=r^2-2r^3+r^4-2r-r^2+4=0$. For the positive weights
$\lambda=(2r-1,r^2-1,1)$, direct expansion gives
\[
\sum_i\lambda_ig_i(z)=(z-p)^{\mathsf T}
\begin{pmatrix}2r&-1\\-1&r^2\end{pmatrix}(z-p),
\]
and this matrix is positive definite, with determinant $2r^3-1=3$. A feasible
point makes the left side nonpositive, so it equals $p$.

The quadratic parts of $R_1,R_2,R_3$ are
$\left(\begin{smallmatrix}5&-1\\-1&2\end{smallmatrix}\right)$,
$\left(\begin{smallmatrix}2&-1\\-1&5\end{smallmatrix}\right)$ and
$\left(\begin{smallmatrix}5&-4\\-4&5\end{smallmatrix}\right)$, each positive
definite with determinant nine. The transformation $3I+\mathbf1\mathbf1^{\mathsf
T}$ from $(g_i)$ to $(R_i)$ is invertible, so the span stays three. With
$\bar\lambda=\sum_i\lambda_i$, the weights $\omega_i=(\lambda_i-\bar\lambda/6)/3$ satisfy
$\sum_i\omega_iR_i=\sum_i\lambda_ig_i$, and they are positive:
$18\omega_1=10r-r^2-5$, $18\omega_2=5r^2-2r-5$ and $18\omega_3=7-2r-r^2$,
and $5/4<r<4/3$ gives the lower bounds $95/16$, $5/16$ and $23/9$.

Finally, a nonnegative aggregate with rational weights vanishes at $p$, so its
gradient vanishes there. Its gradient at $p$ is a rational combination of
$1,r,r^2$, which are linearly independent over $\Q$; hence all its quadratic
and linear coefficients vanish. A nonnegative combination of nonzero positive
semidefinite Hessians is zero only if all weights are zero, since traces are
positive.
\end{proof}

Thus each set $R_i\le0$ is a full-dimensional compact ellipsoid with rational
data, and three of them can meet in a single irrational point. The point
$(\sqrt[3]2,\sqrt[3]4)$ is also the classical example of an irrational
singleton cut out by a convex cubic inequality~\cite{BienstockDelPiaHildebrand2023};
the native PSD quadratic presentation is what matters here. The last
statement explains why the zero-margin step of Theorem~\ref{thm:qc-span-two}
cannot be copied at span three: an exposing aggregate with minimum zero may
need irrational weights. It does not exclude rational infeasibility
certificates with a positive margin, which are proved to exist in
Section~\ref{app:qc-infeasible}. Every coordinate of a rational convex
singleton has exactly one real conjugate, and every such number occurs
(Theorem~\ref{thm:singleton-field}); three rational ellipsoids already
realize every such number of degree at least three
(Corollary~\ref{cor:algebraic-three-quadrics}).

\subsection{Singleton spectrahedra}
\label{app:qc-spectra}

Rational semidefinite constraints are a second natural way to pin down an
algebraic number, and they reach different numbers than convex polynomial
inequalities. A rational convex singleton has a coordinate field with exactly
one real embedding (Theorem~\ref{thm:singleton-field}). A rational affine
semidefinite constraint in one scalar variable reaches exactly the totally
real numbers, and its matrix size is at least twice the degree. If the unique
feasible matrix has corank one, the coordinate field again has exactly one
real embedding. These statements show that the representation, and not only
convexity of the feasible set, decides which algebraic numbers can be imposed
exactly.

\begin{theorem}[One scalar variable]
\label{thm:qc-pencil}
For a real number $\alpha$ the following are equivalent.
\begin{enumerate}[label=(\roman*),leftmargin=*]
\item There are rational symmetric matrices $A$ and $B$ of some size with
$\{x\in\R:A+xB\succeq0\}=\{\alpha\}$.
\item $\alpha$ is algebraic and every root of its minimal polynomial is real.
\end{enumerate}
If these hold and $d=[\Q(\alpha):\Q]$, the least possible size of $A$ and
$B$ is $2d$. A pencil of size $2d$ can be computed in time polynomial in the
length of the minimal polynomial and of a rational interval isolating
$\alpha$ whose endpoints are not roots.
\end{theorem}

\begin{proof}
Assume (i). Let $K_0=\ker A\cap\ker B$, a rational subspace. A rational
congruence with blocks for $K_0^\perp$ and $K_0$ turns the pencil into
$\diag(\hat A+x\hat B,0)$ with the same feasible set. The reduced pencil has
no common kernel vector: if $\hat Av=\hat Bv=0$ for the coordinate vector $v$
of some $y\in K_0^\perp$, then $Ay$ and $By$ lie in the range of $A$ and $B$,
which is contained in $K_0^\perp$, and they are orthogonal to $K_0^\perp$.
Hence $Ay=By=0$, so $y\in K_0\cap K_0^\perp=\{0\}$.

Let $C=\hat A+\alpha\hat B\succeq0$, and suppose that
$(C+(z-\alpha)\hat B)v=0$ for a nonreal $z$ and a complex $v\ne0$. Then
$v^*Cv+(z-\alpha)v^*\hat Bv=0$ with both quadratic forms real, so
$v^*\hat Bv=0$, then $v^*Cv=0$, hence $Cv=0$, $\hat Bv=0$ and $\hat Av=0$.
The real or imaginary part of $v$ would be a nonzero common kernel vector.
Therefore $\mathcal D(x)=\det(\hat A+x\hat B)\in\Q[x]$ is nonzero and has only
real roots. The matrix $C$ is singular, since otherwise it would remain
positive definite near $\alpha$. Thus $\mathcal D(\alpha)=0$, so $\alpha$ is
algebraic and its minimal polynomial, a factor of $\mathcal D$, has only real
roots.

Next, $\alpha$ is a double root of $\mathcal D$. If $C$ has corank at least
two, its adjugate vanishes and $\mathcal D'(\alpha)=\tr(\operatorname{adj}(C)
\hat B)=0$. If $C$ has corank one with unit kernel vector $v$, then in an
orthonormal basis adapted to $v^\perp$ and $v$,
\[
C+t\hat B=\begin{pmatrix}C_1+tE&tb\\ tb^{\mathsf T}&t\beta\end{pmatrix},
\qquad C_1\succ0,\quad\beta=v^{\mathsf T}\hat Bv .
\]
If $\beta\ne0$, then for small $t$ with $t\beta>0$ the block $C_1+tE$ is
positive definite and the Schur complement $t\beta-t^2b^{\mathsf T}
(C_1+tE)^{-1}b$ is positive, so $\alpha+t$ is feasible. Hence $\beta=0$, and
$\mathcal D'(\alpha)$, which is a multiple of $\beta$, vanishes. The minimal
polynomial $\mu$ of $\alpha$ divides $\mathcal D$ and $\mathcal D'$; since
$\mu$ is irreducible and does not divide $\mu'$, $\mu^2$ divides $\mathcal D$.
Hence $2d\le\deg\mathcal D\le$ size of the pencil.

Assume (ii), and let $\alpha_1<\cdots<\alpha_d$ be the conjugates, with
$\alpha=\alpha_\kappa$. In the power basis $1,\alpha,\ldots,\alpha^{d-1}$ of
$\Q(\alpha)$, let $M$ be the rational matrix of multiplication by $\alpha$ and
$G_{ij}=\tr(M^{i+j})$ the trace form. With the real Vandermonde matrix
$V_{ji}=\alpha_j^i$, one has $G=V^{\mathsf T}V\succ0$ and
$VM=\diag(\alpha_1,\ldots,\alpha_d)V$. For rational $\varrho$ that is not a
conjugate, the rational symmetric affine pencil
\[
\mathcal L_\varrho(x)=G(M-xI)(M-\varrho I)^{-1}
=V^{\mathsf T}\diag\Bigl(\frac{\alpha_j-x}{\alpha_j-\varrho}\Bigr)V
\]
is congruent to the displayed diagonal matrix. Choose rational
$\varrho_-<\alpha<\varrho_+$ such that $[\varrho_-,\varrho_+]$ contains no other
conjugate. The $\kappa$th diagonal entry of $\mathcal L_{\varrho_-}(x)$ is
nonnegative exactly when $x\le\alpha$, and that of $\mathcal L_{\varrho_+}(x)$
exactly when $x\ge\alpha$. At $x=\alpha$ every other entry is positive,
because $\alpha_j-\alpha$, $\alpha_j-\varrho_-$ and $\alpha_j-\varrho_+$ have
the same sign for $j\ne\kappa$. Hence
$\diag(\mathcal L_{\varrho_-}(x),\mathcal L_{\varrho_+}(x))\succeq0$ exactly
when $x=\alpha$, at size $2d$. The matrices are obtained from the companion
matrix by $2d-2$ powers, traces, and two inversions, all with polynomial
coefficient bit length.
\end{proof}

That positive semidefinite pencils, including singular ones, have only real
finite eigenvalues is classical~\cite{LiangLiBai2013}; the proof above is
included to make the arithmetic consequence and the singular case explicit.
The unique feasible matrix of the construction has corank two, and the
determinant of the pencil is
\[
\det\bigl(\diag(\mathcal L_{\varrho_-}(x),\mathcal L_{\varrho_+}(x))\bigr)
=\frac{\det(G)^2\,\mu(x)^2}{\mu(\varrho_-)\,\mu(\varrho_+)} .
\]
For
example $\sqrt2$ is imposed by a pencil of size four, while $\sqrt[3]2$, which
has two nonreal conjugates, cannot be imposed in one variable by any rational
pencil. Auxiliary scalar variables change the statement, as the next theorem
shows.

\begin{theorem}[Corank one]
\label{thm:qc-corank-one}
Let $A(x)=A_0+\sum_{i=1}^nx_iA_i$ with rational symmetric $A_i$, and suppose
$\{x\in\R^n:A(x)\succeq0\}=\{x^*\}$ and $A(x^*)$ has corank one. Then
$K=\Q(x_1^*,\ldots,x_n^*)$ is a number field with exactly one real embedding,
and so is every subfield; in particular every coordinate has exactly one real
conjugate. Conversely, every number field with exactly one real embedding is
the coordinate field of such a singleton with $n\le2$.
\end{theorem}

\begin{proof}
The singleton is defined over $\Q$ by sign conditions on principal minors.
The real algebraic numbers form a real closed field, and by the
Tarski--Seidenberg transfer principle the existential statement that the
pencil has a positive semidefinite value holds over that
field~\cite{BasuPollackRoy2006}. Uniqueness gives $x^*$ algebraic coordinates,
so $K$ is a number field.

Let $R=A(x^*)$. It has entries in $K$ and rank $m-1$, so its kernel is spanned
by some $v\in K^m$; normalize a coordinate $v_\iota=1$. If
$v^{\mathsf T}A_iv\ne0$ for some $i\ge1$, choose $\pm e_i$ so that the
direction $D$ satisfies $v^{\mathsf T}Dv>0$; the Schur-complement argument of
Theorem~\ref{thm:qc-pencil} makes $R+tD$ positive definite for small $t>0$,
a second feasible point. Hence $v^{\mathsf T}A_iv=0$ for $i\ge1$, and then
$v^{\mathsf T}A_0v=v^{\mathsf T}Rv=0$. Next, the linear system $A(x)v=0$ in
$x$ has the unique solution $x^*$: if $D=\sum_i\delta_iA_i$ with $\delta\ne0$
satisfied $Dv=0$, symmetry would make $D$ block diagonal in the decomposition
$v^\perp\oplus\R v$, and $R+tD$ would stay positive semidefinite for all
small $\abs t$. The coefficients of this system lie in $F=\Q(v_1,\ldots,v_m)$,
so Gaussian elimination gives $x^*\in F^n$ and $K=F$.

Let $\sigma$ be a real embedding of $F$ and $u=\sigma(v)\in\R^m$. Since the
$A_i$ are rational, $u^{\mathsf T}A_iu=\sigma(v^{\mathsf T}A_iv)=0$ for
$0\le i\le n$, and therefore $u^{\mathsf T}Ru=u^{\mathsf T}A_0u+
\sum_ix_i^*u^{\mathsf T}A_iu=0$, where $R$ is the original matrix, not a
conjugate. Positive semidefiniteness gives $Ru=0$, so $u\in\R v$, and
$u_\iota=1=v_\iota$ gives $u=v$. Hence $\sigma$ fixes the generators of $F$:
the field has exactly one real embedding. Its degree is then odd. For a
subfield $F'$, the degree $[F:F']$ is odd; a real embedding of $F'$ extends to
$F$ by sending a primitive element of $F/F'$ to a real root of the image of
its odd-degree minimal polynomial. Since $F$ has one real embedding, so does
$F'$.

For the converse, the pencil $\left(\begin{smallmatrix}1&x\\x&0\end{smallmatrix}
\right)$ gives $\{0\}$ with corank one, which covers $\Q$. Let $K\ne\Q$ have
exactly one real embedding, let $\alpha\in K\subset\R$ be a primitive element,
and let $\mu$ be its monic minimal polynomial, of degree $d\ge3$ with exactly
one real root. Let $\mathcal C$ be the companion matrix of $\mu$ with ones on
the superdiagonal and last row $(-\mu_0,\ldots,-\mu_{d-1})$, where
$\mu=T^d+\sum_{i<d}\mu_iT^i$, and let
$\nu(T)=(1,T,\ldots,T^{d-1})^{\mathsf T}$, so that $\mathcal C\nu(\gamma)=
\gamma\nu(\gamma)$ for every root $\gamma$. Let $U$ be the real span of the
real and imaginary parts of $\nu(\gamma)$ over the nonreal roots; it is
$\mathcal C$-invariant, $\R^d=\R\nu(\alpha)\oplus U$, and $\mathcal C-\alpha I$
maps $\R^d$ onto $U$ with kernel $\R\nu(\alpha)$. The rational linear space
$\mathcal L=\{B=B^{\mathsf T}:\nu(T)^{\mathsf T}B\nu(T)\equiv0\bmod\mu(T)\}$
contains the real matrix whose form in the basis
$\nu(\alpha),\Re\nu(\gamma_1),\Im\nu(\gamma_1),\ldots$ is
$\diag(0,1,1,\ldots,1)$: its form vanishes at every root, because
$\nu(\gamma)^{\mathsf T}B\nu(\gamma)=1-1+2\mathrm i\cdot0$. By density of
rational points in $\mathcal L$ and openness of positive definiteness on $U$,
there is a rational $B\in\mathcal L$ positive definite on $U$. Put
\[
Q(s,t)=\mathcal C^{\mathsf T}B\mathcal C-s(\mathcal C^{\mathsf T}B+B\mathcal C)+tB .
\]
Then $Q(\alpha,\alpha^2)=(\mathcal C-\alpha I)^{\mathsf T}B(\mathcal C-\alpha I)$
is positive semidefinite with kernel $\R\nu(\alpha)$, so it has corank one.
Every $Q(s,t)$ satisfies $\nu(\alpha)^{\mathsf T}Q(s,t)\nu(\alpha)=0$, because
$\nu(\alpha)^{\mathsf T}B\nu(\alpha)=0$. If $Q(s,t)\succeq0$, it therefore
annihilates $\nu(\alpha)$, which with $w=B\nu(\alpha)$ reads
\[
(\alpha-s)\mathcal C^{\mathsf T}w+(t-s\alpha)w=0 .
\]
Here $w\ne0$: otherwise every rational row of $B$ would vanish on the power
basis of $\alpha$, forcing $B=0$. The vectors $w$ and $\mathcal C^{\mathsf T}w$
are independent: an eigenvector relation $\mathcal C^{\mathsf T}w=\varkappa w$
with real $\varkappa$ forces $\varkappa=\alpha$, then
$(\gamma-\alpha)w^{\mathsf T}\nu(\gamma)=0$ for every root $\gamma$, and
together with $w^{\mathsf T}\nu(\alpha)=0$ this gives $w=0$. Hence
$(s,t)=(\alpha,\alpha^2)$ is the unique feasible point, and its coordinate
field is $\Q(\alpha)=K$.
\end{proof}

Combined with Theorem~\ref{thm:qc-pencil}, a one-variable singleton whose
feasible matrix has corank one is rational, since its field would be both
totally real and have only one real embedding. The size-$2d$ pencils of
Theorem~\ref{thm:qc-pencil} have corank two, consistent with this. The
statement does not apply to higher corank, to non-singleton feasible sets, or
to projections of spectrahedra.

\subsection{Short inputs with long expanded outputs}
\label{app:qc-short}

Theorem~\ref{thm:qc-sharp} attains the degree bound but says nothing about
the size of the instances, and a large degree has consequences for output
length only if the input stays short. This subsection constructs explicit
short instances whose value and optimizer have degree $n^{\Omega(k)}$. It then
makes the minimal polynomials of the value and of one input coordinate dense,
so that sparse coefficient lists are long as well. The consequence is
unconditional: no algorithm whose output must contain such a list can run in
time $\phi(k)L^C$. Neither decision nor approximation is affected, and
the instances have short descriptions as sums of block numbers.

Each block is a trust-region problem. Its secular equation and its
reduction to a polynomial of degree $2r$ are standard~\cite{AdachiIwataNakatsukasaTakeda2017};
the new point is that this polynomial is Eisenstein for short explicit data,
and that congruences make the fields of different blocks independent.

\begin{proposition}[One block]
\label{prop:qc-block}
Let $\ell\equiv1\pmod4$ be prime, $r=\ell-1$, $\kappa\in\{1,\ldots,\ell-1\}$
with $\kappa^2\equiv-1\pmod\ell$, and $\gamma$ a positive integer with
$\gamma\equiv\kappa\pmod{\ell^2}$. The problem
\begin{equation}
\label{eq:qc-block}
\min\Bigl\{\tfrac12\sum_{j=1}^rjx_j^2-\gamma\sum_{j=1}^rjx_j:\
\sum_{j=1}^rx_j^2\le1\Bigr\}
\end{equation}
has the unique optimizer $x_j^*=\gamma j/(j+\lambda)$, where $\lambda>1$ is
the unique positive solution of $\sum_j\gamma^2j^2/(j+\lambda)^2=1$. Let
$\mathcal F(T)=\prod_{j=1}^r(T+j)$, $\mathcal F_j=\mathcal F/(T+j)$, and
\[
\mathcal S(T)=\mathcal F(T)^2-\gamma^2\sum_{j=1}^rj^2\mathcal F_j(T)^2 .
\]
Then $\mathcal S$ is monic of degree $2r$, Eisenstein at $\ell$, and
$\mathcal S(\lambda)=0$. Each coordinate $x_j^*$ and the optimal value
$\theta$ generate the field $\Q(\lambda)$, which has degree $2r$.
\end{proposition}

\begin{proof}
The objective is strictly convex and the ball is compact with interior point
$0$. The unconstrained minimizer $(\gamma,\ldots,\gamma)$ has norm
$\gamma\sqrt r>1$, so the constraint is active, and the
Karush--Kuhn--Tucker conditions give $x_j=\gamma j/(j+\lambda)$ with
$\sum_jx_j^2=1$. The left side of the secular equation decreases strictly on
$[0,\infty)$ from $r\gamma^2>1$ to $0$, so $\lambda$ is unique. Since
$\kappa\ne1$ (as $1\not\equiv-1$), $\gamma\ge2$, and the left side at
$\lambda=1$ is at least $r\gamma^2/4\ge4$; hence $\lambda>1$. Multiplying the
secular equation by $\mathcal F(\lambda)^2$ gives $\mathcal S(\lambda)=0$.

Modulo $\ell$, $\mathcal F\equiv T^r-1$ and $\gamma^2\equiv-1$, so
$\mathcal S\equiv\mathcal F^2+\sum_jj^2\mathcal F_j^2$. Let $a=-j$ be a root of
$T^r-1$ in $\mathbb F_\ell$. Then $\mathcal F_i(a)=0$ for $i\ne j$, and
$\mathcal F_j(a)=\mathcal F'(a)=ra^{r-1}=-a^{-1}$, so $\mathcal S(a)=
j^2a^{-2}=1=a^{2r}$. From $\mathcal F''=2\mathcal F_j'+(T+j)\mathcal F_j''$ we
get $\mathcal F_j'(a)/\mathcal F_j(a)=\mathcal F''(a)/(2\mathcal F'(a))=
(r-1)/(2a)=r/a$ in $\mathbb F_\ell$, because $r\equiv-1$. Hence
$\mathcal S'(a)=2j^2\mathcal F_j(a)^2\,r/a=2ra^{2r-1}$. The difference
$\mathcal S-T^{2r}$ has degree less than $2r$ and vanishes to second order at
the $r$ distinct roots of $T^r-1$, so $\mathcal S\equiv T^{2r}\pmod\ell$. The
constant term is $\mathcal S(0)=(r!)^2(1-r\gamma^2)$. Write $\kappa^2+1=m\ell$;
then $1\le m\le\ell-2$ and
$(1-r\kappa^2)/\ell=1+m-m\ell\equiv1+m\not\equiv0\pmod\ell$. Since
$\gamma^2\equiv\kappa^2\pmod{\ell^2}$, the constant term has $\ell$-adic
valuation exactly one. Thus $\mathcal S$ is Eisenstein, hence irreducible,
and $[\Q(\lambda):\Q]=2r$~\cite{Neukirch1999}.

Every $x_j^*$ is nonzero, lies in $\Q(\lambda)$, and recovers
$\lambda=\gamma j/x_j^*-j$. For the value, the active constraint and the
identity $j^3/(j+\lambda)^2=j^2/(j+\lambda)-\lambda j^2/(j+\lambda)^2$ give
\begin{equation}
\label{eq:qc-block-value}
2\theta=-\lambda-\gamma^2\sum_{j=1}^r\frac{j^2}{j+\lambda}\in\Q(\lambda).
\end{equation}
Extend the $\ell$-adic valuation $v_\ell$, normalized by $v_\ell(\ell)=1$, to
a finite extension of $\Q_\ell$ containing $\lambda$. Since $\mathcal S$ is
Eisenstein, $v_\ell(\lambda)=1/(2r)$, and the ramification index of a local
field bounds the denominators of its valuations~\cite{Neukirch1999}. The
integers $j\le r$ are units, so
\[
2\theta=-\gamma^2\sum_jj+(r\gamma^2-1)\lambda
+\sum_{e\ge2}(-1)^{e+1}\gamma^2\Bigl(\sum_jj^{1-e}\Bigr)\lambda^e
\]
converges. The constant term $-\gamma^2\ell(\ell-1)/2$ has valuation one, and
the coefficient of $\lambda$ is divisible by $\ell$. For $2\le e\le r$, the
power sum $\sum_{a\in\mathbb F_\ell^\times}a^{1-e}$ vanishes because
$\ell-1\nmid e-1$, so these coefficients are divisible by $\ell$. For $e=r+1$
the power sum is $r\equiv-1$, a unit. Later coefficients are integral. Hence
the term $e=r+1$ has the unique least valuation and
$v_\ell(2\theta)=(r+1)/(2r)$. Since $r$ is even, $\gcd(r+1,2r)=1$, so the
ramification index of $\Q_\ell(\theta)$ over $\Q_\ell$ is at least $2r$, and
$[\Q(\theta):\Q]\ge2r$. With $\theta\in\Q(\lambda)$, this gives
$\Q(\theta)=\Q(\lambda)$.
\end{proof}

\begin{lemma}[Splitting at a larger prime]
\label{lem:qc-split}
In Proposition~\ref{prop:qc-block}, drop the congruence condition on
$\gamma$. If $q>r$ is an odd prime dividing $\gamma$, then $\mathcal S$ has
$2r$ distinct roots in $\Q_q$.
\end{lemma}

\begin{proof}
Fix $i$ and substitute $T=-i+\gamma iu$ in the secular equation. The term
$j=i$ becomes $u^{-2}$, and multiplying by $u^2\prod_{j\ne i}(j-i+\gamma iu)^2$
gives an integer polynomial equation that reduces modulo $q$ to
$(u^2-1)\prod_{j\ne i}(j-i)^2=0$, with nonzero product because
$0<\abs{j-i}<q$. The roots $u=\pm1$ are simple modulo $q$, so Hensel's
lemma~\cite{Neukirch1999} lifts them to $u_\pm\in\Z_q$. The factors $j-i+\gamma iu_\pm$ are units and
$u_\pm\ne0$, so $T_{i,\pm}=-i+\gamma iu_\pm$ are roots of $\mathcal S$. They
are distinct for the two signs, and they are congruent to $-i$ modulo $q$, so
roots for different $i$ are distinct.
\end{proof}

The reduction of $\mathcal S$ modulo $q$ is $\mathcal F^2$, which is not
squarefree; the rescaling around each pole is needed.

\begin{theorem}[Products of blocks]
\label{thm:qc-blocks}
Let $\ell_1<\cdots<\ell_k$ be primes congruent to one modulo four, choose
$\kappa_i$ as in Proposition~\ref{prop:qc-block}, put
$M_i=\prod_{i'>i}\ell_{i'}$, choose $1\le\varrho_i<\ell_i^2$ with
$M_i\varrho_i\equiv1\pmod{\ell_i^2}$, and let $\gamma_i=\kappa_iM_i\varrho_i$.
Let $\lambda_i$, $\theta_i$ and $E_i=\Q(\lambda_i)$ belong to block $i$, and put
$r_i=\ell_i-1$, $n=\sum_ir_i$ and $\mathcal D=\prod_i2r_i$.
\begin{enumerate}[label=(\alph*),leftmargin=*]
\item $[E_1\cdots E_k:\Q]=\mathcal D$.
\item For positive rational weights $\omega_i$, the problem of minimizing
$\sum_i\omega_i\varphi_i(x_i)$ subject to $\norm{x_i}^2\le1$, where
$\varphi_i$ is the objective of~\eqref{eq:qc-block} for block $i$, has the
unique optimizer formed by the block optimizers; its coordinates generate
$E_1\cdots E_k$, and its Hessian span is exactly $k$.
\item Some integer $1\le t\le1+(k-1)\mathcal D(\mathcal D-1)/2$ makes the
optimal value for the weights $\omega_i=t^{i-1}$ generate $E_1\cdots E_k$,
so it has degree $\mathcal D$.
\item Let the weights satisfy $1\le\omega_1\le\cdots\le\omega_k$, as those
of (c) do, and put $\bar\gamma=\max_i\gamma_i\ell_i^2$ and
$\eta=1/(2k\omega_k\bar\gamma)$. Replacing the rows by
\[
\norm{x_i}^2+\eta\sum_{i'\ne i}\norm{x_{i'}}^2-1-\eta(k-1)\le0\qquad(1\le i\le k)
\]
keeps the optimizer and the value, makes every constraint Hessian positive
definite, and keeps the span equal to $k$. The feasible set is compact and
$0$ is strictly feasible.
\item With the weights of (c) and the rows of (d), every coefficient has
$O(k^2\log\ell_k)$ bits, and the input length is
$L=O\bigl((k+1)n^2(1+k^2\log\ell_k)\bigr)$ in the native encoding and
$L=O\bigl((k+1)n^2(n+k^2\log\ell_k)\bigr)$ in the explicit format.
\end{enumerate}
\end{theorem}

\begin{proof}
(a) By construction $\gamma_i\equiv\kappa_i\pmod{\ell_i^2}$, so each
$\mathcal S_i$ is Eisenstein at $\ell_i$, and $\ell_i$ divides $\gamma_{i'}$
for $i'<i$. Suppose $[E_1\cdots E_{i-1}:\Q]=\prod_{i'<i}2r_{i'}$. For
$i'<i$, $\ell_i>r_{i'}$ and Lemma~\ref{lem:qc-split} shows that
$\mathcal S_{i'}$ splits in $\Q_{\ell_i}$. The splitting field of
$\prod_{i'<i}\mathcal S_{i'}$ is generated by roots that all lie in
$\Q_{\ell_i}$, so it embeds in $\Q_{\ell_i}$, and so does $E_1\cdots E_{i-1}$.
The Eisenstein polynomial $\mathcal S_i$ is irreducible over $\Q_{\ell_i}$,
hence over the image of $E_1\cdots E_{i-1}$, and adjoining $\lambda_i$
multiplies the degree by $2r_i$. Pairwise trivial intersections of the
fields would not suffice for this conclusion.

(b) The problem is separable with positive weights, so its optimizer consists
of the block optimizers. By Proposition~\ref{prop:qc-block} the coordinates
of block $i$ generate $E_i$. The constraint Hessians have disjoint nonempty
diagonal supports, hence are independent.

(c) The values $\theta_i$ generate $E=E_1\cdots E_k$. Two distinct
embeddings $\sigma,\tau$ of $E$ into $\C$ differ on some $\theta_i$, so
$\sum_iT^{i-1}(\sigma(\theta_i)-\tau(\theta_i))$ is a nonzero polynomial of degree at
most $k-1$. There are $\mathcal D(\mathcal D-1)/2$ pairs, so some integer in
the stated range is a root of none of these polynomials, and
$\sum_it^{i-1}\theta_i$ then has $\mathcal D$ distinct conjugates.

(d) Each Hessian is $2\diag(\eta,\ldots,1,\ldots,\eta)\succ0$, and the rows
arise from the ball rows by the invertible matrix
$N=(1-\eta)I+\eta\mathbf1\mathbf1^{\mathsf T}$, so the span stays $k$. At the
old optimizer every new row is active. From the secular equation,
$\lambda_i<\gamma_i(\sum_jj^2)^{1/2}<\bar\gamma$, and $\lambda_i>1$ by
Proposition~\ref{prop:qc-block}. Put $a_i=\omega_i\lambda_i$. Then $a_i>1$
because $\omega_i\ge1$, and $\Sigma=\sum_ia_i<k\omega_k\bar\gamma$ because
$\omega_i\le\omega_k$. Stationarity of
$\sum_i\omega_i\varphi_i+\tfrac12\sum_i\mu_i(\text{row }i)$ at the old
optimizer requires $N\mu=a$, since every block optimizer is nonzero. The
solution is
\[
\mu_i=\frac{a_i-\eta\Sigma/(1+\eta(k-1))}{1-\eta}>0,
\]
because $\eta\Sigma<1/2<a_i$. The Karush--Kuhn--Tucker conditions are
sufficient for this convex problem, and strict convexity of the objective
gives uniqueness. The feasible set changes, but the optimizer does not.

The conditions on the weights cannot be dropped. For $k=2$ and $\omega_2=1$,
the number $\eta=1/(4\bar\gamma)$ does not depend on $\omega_1$, and the
unique solution of $N\mu=a$ has
$\mu_1=(\omega_1\lambda_1-\eta\lambda_2)/(1-\eta^2)$ and
$\mu_2=(\lambda_2-\eta\omega_1\lambda_1)/(1-\eta^2)$. Thus $\mu_2<0$ for
large $\omega_1$, and $\mu_1<0$ for small $\omega_1>0$. Since $0$ is a
Slater point, the Karush--Kuhn--Tucker conditions are necessary, so in both
cases the old optimizer is not optimal for the new rows.

(e) We have $\log\gamma_i=O(k\log\ell_k)$, $\log\mathcal D=O(k\log\ell_k)$,
$\log t^{k-1}=O(k^2\log\ell_k)$, and $\log(1/\eta)=O(k^2\log\ell_k)$. There are
$k+1$ quadratics in $n$ variables, each with $O(n^2)$ coefficients; in the
explicit format each monomial also carries an exponent vector of $O(n)$
bits.
\end{proof}

\begin{corollary}[No parameterized bound for dense output]
\label{cor:qc-no-fpt-output}
For every function $\phi$ and constant $C$, no algorithm that returns the
dense coefficient list of the minimal polynomial of the optimal value of
every rational convex quadratically constrained problem can run in time
$\phi(k)L^C$, where $k$ is the Hessian span.
\end{corollary}

\begin{proof}
Fix an integer $k>3C$. By the prime number theorem for arithmetic
progressions~\cite{Davenport2000}, every interval $[R,2R]$ with $R$ large
contains at least $k$ primes congruent to one modulo four. Building
Theorem~\ref{thm:qc-blocks}(c)--(e) from such primes gives $n<2kR$, input
length $L=O_k(R^2\log R)$ in the native encoding and $L=O_k(R^3)$ in the
explicit format, and a value of degree $\mathcal D\ge(2(R-1))^k$, whose
dense list has $\mathcal D+1$ entries. For large $R$ this exceeds
$\phi(k)L^C$ in both formats.
\end{proof}

The weight in Theorem~\ref{thm:qc-blocks}(c) is shown to exist; the following
lemma gives a larger weight that can be written down directly, which makes
the whole family constructible in time polynomial in $n$.

\begin{lemma}[An explicit primitive weight]
\label{lem:qc-weight}
Let $\beta_1,\ldots,\beta_k$ be real algebraic numbers generating a field
$E$, and let each $\beta_i$ be a root of a nonzero integer polynomial of
degree at most $d$ whose coefficients have absolute values summing to at most
$2^H$, $H\ge1$. Then for $t=1+2^{(H+2)d^2+H+3}$ the number
$\sum_it^{i-1}\beta_i$ generates $E$.
\end{lemma}

\begin{proof}
All conjugates of $\beta_i$ are roots of the given polynomial, so they have
modulus at most $1+2^H\le2^{H+1}$. Two distinct conjugates of $\beta_i$ are
roots of its primitive minimal polynomial, whose degree is $e\le d$ and whose
leading coefficient has modulus at most $2^H$. Its discriminant is a nonzero
integer equal to the leading coefficient to the power $2e-2$ times the
product of the squared root differences. Bounding all other factors gives a
root separation of at least $\delta_\star=2^{-(H+2)d^2}$. Let $\sigma\ne\tau$
be embeddings of $E$ and $j$ the largest index with
$\sigma(\beta_j)\ne\tau(\beta_j)$. Then
\[
\Bigl\lvert\sum_it^{i-1}(\sigma\beta_i-\tau\beta_i)\Bigr\rvert
\ge t^{j-1}\delta_\star-2^{H+2}\frac{t^{j-1}}{t-1}>0,
\]
because $2^{H+2}/(t-1)=2^{-(H+2)d^2-1}<\delta_\star$.
\end{proof}

For block $i$, $2\theta_i$ is a root of the resultant
$\operatorname{Res}_T\bigl(\mathcal S_i(T),\,Y\mathcal F_i(T)+T\mathcal F_i(T)
+\gamma_i^2\sum_jj^2\mathcal F_{i,j}(T)\bigr)$, by~\eqref{eq:qc-block-value}.
This resultant has degree $2r_i$ in $Y$, since $\mathcal S_i$ is monic and
$\mathcal F_i$ does not vanish at its roots. Its Sylvester matrix has size
$3r_i+1$ and entries whose coefficient sums are at most
$((r_i+1)!)^2(2+\gamma_i^2r_i^3)$, which gives
$H=O(\ell_k^2\log\ell_k+k\ell_k\log\ell_k)$ and $d=2\ell_k$ in
Lemma~\ref{lem:qc-weight}. The resulting weights have
$O(k\ell_k^4\log\ell_k+k^2\ell_k^3\log\ell_k)$ bits, and the analogue of
Corollary~\ref{cor:qc-no-fpt-output} holds with $L=O_k(R^6\log R)$ in either
format and $k>6C$.

That the sparsity of a polynomial depends on the basis, in particular on a
shift of the variable, is well known~\cite{GiesbrechtRoche2010}; the next
lemma only chooses a short shift whose ordinary expansion is dense.

\begin{lemma}[Translation makes coefficients nonzero]
\label{lem:qc-translate}
Let $\Pi(X)=\sum_{i=0}^{\mathcal D}\pi_iX^i\in\Q[X]$ with
$\pi_{\mathcal D}\ne0$. Some integer $0\le s\le\mathcal D(\mathcal D+1)/2$
makes every coefficient of $\Pi(X-s)$ nonzero. If $\Pi$ is the minimal
polynomial of $\vartheta$, then $\Pi(X-s)$ is that of $\vartheta+s$, and
integer translation preserves primitivity.
\end{lemma}

\begin{proof}
The coefficient of $X^j$ in $\Pi(X-s)$ is
$\sum_{i\ge j}\pi_i\binom ij(-s)^{i-j}$, a polynomial in $s$ of degree
$\mathcal D-j$ with leading coefficient $\pi_{\mathcal D}\binom{\mathcal
D}j(-1)^{\mathcal D-j}\ne0$. Together these polynomials have at most
$\sum_j(\mathcal D-j)=\mathcal D(\mathcal D+1)/2$ integer roots. The
substitution $X\mapsto X-s$ is an automorphism of $\Q[X]$ and of $\Z[X]$.
\end{proof}

\begin{theorem}[Sparse lists]
\label{thm:qc-sparse}
The instances of Theorem~\ref{thm:qc-blocks}(c)--(e) can be modified, keeping
positive definite constraint Hessians, Hessian span exactly $k$, a compact
strictly feasible set, a unique optimizer, and the input lengths of
Theorem~\ref{thm:qc-blocks}(e) up to constant factors, so that the optimal
value and the first
coordinate of the optimizer both have degree $\mathcal D$ and minimal
polynomials with exactly $\mathcal D+1$ nonzero coefficients. Consequently no
algorithm can run in time $\phi(k)L^C$ if it must return either the explicit
list of nonzero coefficients of the minimal polynomial of the optimal value,
or the minimal polynomials of all coordinates of the optimizer.
\end{theorem}

\begin{proof}
Add to the objective an integer constant $s_0$ from
Lemma~\ref{lem:qc-translate} for the minimal polynomial of the value. For the
coordinate, let $a_i=x_{i,1}^*$ be the first coordinate of block $i$; it
generates $E_i$ because $\lambda_i=\gamma_i/a_i-1$. As in
Theorem~\ref{thm:qc-blocks}(c), some integer
$1\le\kappa'\le1+(k-1)\mathcal D(\mathcal D-1)/2$ makes
$\alpha'=\sum_i\kappa'^{\,i-1}a_i$ primitive for $E$, and
Lemma~\ref{lem:qc-translate} gives an integer $s_1$ for its minimal
polynomial. The affine change of variables
$y_1=x_{1,1}+\sum_{i\ge2}\kappa'^{\,i-1}x_{i,1}+s_1$, with all other
coordinates unchanged, has a unimodular linear part. It maps the optimizer to
a point whose first coordinate is $\alpha'+s_1$. A congruence by an invertible
matrix preserves positive definiteness of every Hessian and the dimension of
their span, and an affine bijection preserves compactness, strict
feasibility, and uniqueness. Each new coefficient is a sum of at most four
products of old coefficients with $\kappa'^{\,i-1}$ and $s_1$, which have
$O(k^2\log\ell_k)$ bits. The output bound follows as in
Corollary~\ref{cor:qc-no-fpt-output}: both requested lists have
$\mathcal D+1$ entries.
\end{proof}

The theorem concerns the ordinary monomial basis of the minimal polynomial of
the requested scalar. It does not exclude a polynomial in a shifted basis, an
annihilating polynomial that is not minimal, another primitive element with
coordinate maps, a tower of fields, the implicit block description given by
the secular equations, or a circuit with root-selection gates; a circuit over
the rationals alone cannot represent these irrational numbers. The shear is a
rational affine map and is undone by one. In the unsheared instances each
coordinate has degree $2r_i$ only, so separate coordinate descriptions are
short there; the shear removes exactly this saving for one input
coordinate. The integers $s_0$, $\kappa'$ and $s_1$ are
shown to exist, which suffices for an output lower bound. Nothing here
concerns the cost of deciding feasibility, comparing values, or approximating
the optimizer.

\paragraph{Feasible points.}
The same output phenomenon occurs for feasibility alone, with ellipsoids and
no objective. The blocks are binomial singletons cut out by three ellipsoids
from the pencil of Theorem~\ref{thm:qc-corank-one}, made explicit with
coefficients of logarithmic length; a Kummer argument controls the joint field.

\begin{lemma}[Independent radicals]
\label{lem:qc-kummer}
Let $d$ be a prime, let $\varpi_1,\ldots,\varpi_k$ be distinct primes, and let
$\alpha_i=\varpi_i^{1/d}$ be the real roots. Then
$[\Q(\alpha_1,\ldots,\alpha_k):\Q]=d^k$, and $\sum_ic_i\alpha_i$ generates this
field for all $c_1,\ldots,c_k\in\Q^\times$.
\end{lemma}

\begin{proof}
Let $\zeta$ be a primitive $d$th root of unity, $\mathbb K=\Q(\zeta)$, of
degree $d-1$, and $\mathbb L=\mathbb K(\alpha_1,\ldots,\alpha_k)$, the
splitting field of $\prod_i(X^d-\varpi_i)$ over $\mathbb K$. An automorphism
of $\mathbb L/\mathbb K$ is determined by the vector $b\in\mathbb F_d^k$ with
$\alpha_i\mapsto\zeta^{b_i}\alpha_i$, and this embeds the Galois group as a
subspace of $\mathbb F_d^k$. If the subspace were proper, some nonzero
$\mathfrak e\in\{0,\ldots,d-1\}^k$ would satisfy $\mathfrak e\cdot b\equiv0$
on it, and $\beta=\prod_i\alpha_i^{\mathfrak e_i}$ would be fixed, hence in
$\mathbb K$. Then $\beta^d=\prod_i\varpi_i^{\mathfrak e_i}$, so the norm of
$\beta$ from $\mathbb K$ to $\Q$ is a rational number $\nu$ with
$\nu^d=\prod_i\varpi_i^{(d-1)\mathfrak e_i}$. Comparing $\varpi_i$-adic
valuations gives $d\mid(d-1)\mathfrak e_i$, hence $\mathfrak e=0$. So all $d^k$
rotations occur. Applying the rotation of one $\alpha_j$ to a relation
$\sum_i\xi_i\alpha_i=0$ with $\xi_i\in\mathbb K$ and subtracting gives
$\xi_j(\zeta-1)\alpha_j=0$, so the $\alpha_i$ are linearly independent over
$\mathbb K$. A rotation fixing $s=\sum_ic_i\alpha_i$ therefore has all
$c_i(\zeta^{b_i}-1)=0$ and is the identity, so
$d^k=[\mathbb K(s):\mathbb K]\le[\Q(s):\Q]\le[\Q(\alpha_1,\ldots,\alpha_k):\Q]
\le d^k$.
\end{proof}

\begin{proposition}[Explicit binomial blocks]
\label{prop:qc-binomial-block}
Let $d\ge3$ be odd, $a\ge2$ an integer, and $\alpha=a^{1/d}$. In time
polynomial in $d+\log a$ one can compute rational quadratics $g_1,g_2,g_3$ in
$d-1$ variables, with positive definite Hessians and coefficients whose
numerators and denominators have $O(\log(da))$ bits, such that
\[
\{x:g_1(x)\le0,\ g_2(x)\le0,\ g_3(x)\le0\}=\{(\alpha,\alpha^2,\ldots,\alpha^{d-1})\},
\]
and $\sum_i\tau_ig_i(x)=(1,x)^{\mathsf T}\mathcal R(1,x)$ for weights
$\tau_i>1/4$ with $\sum_i\tau_i=1$ and a positive semidefinite matrix
$\mathcal R$ whose kernel is spanned by $\nu(\alpha)=(1,\alpha,\ldots,
\alpha^{d-1})$. If $T^d-a$ is irreducible, the three Hessians are linearly
independent.
\end{proposition}

\begin{proof}
Use the construction in the proof of Theorem~\ref{thm:qc-corank-one} for
$\mu=T^d-a$, with an explicit choice of $B$. Here $\mathcal Ce_i=e_{i-1}$ for
$i\ge1$ and $\mathcal Ce_0=ae_{d-1}$, so $\norm{\mathcal C}=a$, and the
eigenvectors of $\mathcal C$ are $\nu(\alpha\zeta^j)$ with
$\zeta=e^{2\pi\mathrm i/d}$. Let $\Delta_\alpha=\diag(1,\alpha,\ldots,
\alpha^{d-1})$ and $B_*=\Delta_\alpha^{-1}(I-d^{-1}\mathbf1\mathbf1^{\mathsf
T})\Delta_\alpha^{-1}$. Since $\Delta_\alpha^{-1}\nu(\alpha\zeta^j)=
(\zeta^{ij})_i$, and for $j\ne0$ both $\sum_i\zeta^{ij}$ and $\sum_i\zeta^{2ij}$
vanish because $d$ is odd, the form of $B_*$ vanishes at every
$\nu(\alpha\zeta^j)$, including $j=0$. Hence $B_*\in\mathcal L$. The space $U$
equals $\Delta_\alpha\mathbf1^\perp$, so $u=\Delta_\alpha w$ with
$w\perp\mathbf1$ gives $u^{\mathsf T}B_*u=\norm w^2\ge a^{-2}\norm u^2$,
while $\norm{B_*}\le1$.

Index rows and columns from zero, and for $0\le r<d$ let $E_r$ have entry one
where $i+j=r$, entry $a$ where $i+j=r+d$, and zero elsewhere. Reducing
$\nu(T)^{\mathsf T}B\nu(T)$ modulo $T^d-a$ shows
$\mathcal L=\{B:\langle E_r,B\rangle_F=0,\ 0\le r<d\}$. The $E_r$ have
disjoint supports and squared norms $D_r=r+1+a^2(d-r-1)\le da^2$, so
$\Pi_{\mathcal L}(X)=X-\sum_rD_r^{-1}\langle E_r,X\rangle_FE_r$ is the
orthogonal projection onto $\mathcal L$. Let $M_0$ be a power of two with
$16da^2\le M_0<32da^2$, let $B_0$ be a symmetric matrix on the grid
$M_0^{-1}\Z$ whose entries are within $1/M_0$ of those of $B_*$, computed by
bisection for $a^{1/d}$ and interval arithmetic, and let $B=\Pi_{\mathcal
L}(B_0)$. Then $\norm{B-B_*}\le\norm{B_0-B_*}_F\le d/M_0\le(16a^2)^{-1}$, so
$\norm B\le2$ and $u^{\mathsf T}Bu\ge(2a^2)^{-1}\norm u^2$ on $U$. Each entry
of $B$ is an entry of $B_0$ minus one correction with denominator $M_0D_r$,
so it has $O(\log(da))$ bits.

Let $Q(s,t)=\mathcal C^{\mathsf T}B\mathcal C-s(\mathcal C^{\mathsf
T}B+B\mathcal C)+tB$ and $\mathcal R=Q(\alpha,\alpha^2)=(\mathcal C-\alpha
I)^{\mathsf T}B(\mathcal C-\alpha I)$, which is positive semidefinite with
kernel $\R\nu(\alpha)$. On $H=\{z:z_0=0\}$, the vector $y=(\mathcal C-\alpha
I)z$ satisfies $z_{j+1}=\alpha z_j+y_j$ for $j\le d-2$, which gives
$\norm z\le ad\norm y$ and $z^{\mathsf T}\mathcal Rz\ge\norm z^2/(2a^4d^2)$.
Since $\norm{\mathcal C^{\mathsf T}B\mathcal C}\le2a^2$,
$\norm{\mathcal C^{\mathsf T}B+B\mathcal C}\le4a$ and $\norm B\le2$, the
restriction of $Q(s,t)$ to $H$ is positive definite whenever
$\abs{s-\alpha},\abs{t-\alpha^2}\le3\epsilon_0$, where
$\epsilon_0=1/(100a^5d^2)$. Choose a dyadic $c\in\Q^2$ within $\epsilon_0/16$
of $(\alpha,\alpha^2)$ in each coordinate, and the triangle with vertices
$c+(-\epsilon_0,-\epsilon_0)$, $c+(2\epsilon_0,-\epsilon_0)$ and
$c+(-\epsilon_0,2\epsilon_0)$. With $\delta=(\alpha,\alpha^2)-c$, the
barycentric weights of $(\alpha,\alpha^2)$ are $1-\tau_2-\tau_3$,
$\tau_2=(1+\delta_1/\epsilon_0)/3$ and $\tau_3=(1+\delta_2/\epsilon_0)/3$, all
larger than $1/4$, and every vertex lies within $3\epsilon_0$ of
$(\alpha,\alpha^2)$. For the vertices $(s_i,t_i)$ put
$g_i(x)=(1,x)^{\mathsf T}Q(s_i,t_i)(1,x)$. The Hessian of $g_i$ is twice the
restriction of $Q(s_i,t_i)$ to $H$, hence positive definite, and
$\sum_i\tau_ig_i(x)=(1,x)^{\mathsf T}\mathcal R(1,x)\ge0$ with equality only
where $(1,x)\in\R\nu(\alpha)$, that is, at $x=(\alpha,\ldots,\alpha^{d-1})$.
Multiplication by $\mathcal C$ or $\mathcal C^{\mathsf T}$ permutes entries
and scales one of them by $a$, so every coefficient of $g_i$ is a sum of a
bounded number of products of entries of $B$ with $1$, $a$, $s_i$ and $t_i$,
and has $O(\log(da))$ bits.

For the last statement, a dependence among the three Hessians can be taken
rational. Since the vectors $(1,s_i,t_i)$ are linearly independent, it yields
$\Theta=\theta_0\mathcal C^{\mathsf T}B\mathcal C-\theta_1(\mathcal C^{\mathsf
T}B+B\mathcal C)+\theta_2B$ with rational $\theta\ne0$ and zero restriction to
$H$. Then $x\mapsto(1,x)^{\mathsf T}\Theta(1,x)$ is a rational affine function
that vanishes at $(\alpha,\ldots,\alpha^{d-1})$, because all three matrices
have zero form at $\nu(\alpha)$; as $\alpha$ has degree $d$, all its
coefficients vanish, and $\Theta=0$. For $\beta=\alpha\zeta^j$ with $1\le j<d$,
evaluating $\Theta$ on $\nu(\beta)$ and $\nu(\bar\beta)$ gives
\[
\bigl(\theta_0\alpha^2-2\theta_1\alpha\cos(2\pi j/d)+\theta_2\bigr)\,
\nu(\beta)^{\mathsf T}B\nu(\bar\beta)=0,
\]
and $\nu(\beta)^{\mathsf T}B\nu(\bar\beta)>0$, since it is the sum of the
values of the form of $B$ at the real and imaginary parts of $\nu(\beta)$,
which lie in $U$ and are not both zero. For $d\ge5$ two different cosines
force $\theta_1=0$, and then $\theta_0=\theta_2=0$ because $\alpha^2$ has degree
$d$. For $d=3$ the relation reads $\theta_0\alpha^2+\theta_1\alpha+\theta_2=0$.
In both cases $\theta=0$, a contradiction.
\end{proof}

\begin{corollary}[Feasible points with dense coordinate polynomials]
\label{cor:qc-feasible-output}
Let $d$ be an odd prime, let $\varpi_1<\cdots<\varpi_k$ be primes, and let
$n=k(d-1)$. There is a rational native PSD system in $n$ variables with $3k$
rows, each defining a full-dimensional ellipsoid, Hessian span exactly $3k$,
and a single feasible point, whose first coordinate is
$R_0+\sum_b\varpi_b^{1/d}$ with $R_0=1+\sum_b\varpi_b$. This coordinate has
degree $d^k$, and its minimal polynomial has all $d^k+1$ coefficients nonzero.
The system is computed in time polynomial in $k$, $d$ and $\log\varpi_k$; its
nonconstant coefficients have $O(\log(kd\varpi_k))$ bits and its constant
terms $O(k\log(kd\varpi_k))$ bits, so its input length is
$L=O(k^3d^2\log(kd\varpi_k))$ in the native encoding and
$L=O(k^3d^2[\log(kd\varpi_k)+kd])$ in the explicit format. Consequently, for
every function $\phi$ and constant $C$, no algorithm can return in time
$\phi(h)L^C$ the explicit
coefficient lists of the minimal polynomials of the coordinates of the unique
feasible point of every rational native PSD system with Hessian span $h$.
\end{corollary}

\begin{proof}
For each $b$, apply Proposition~\ref{prop:qc-binomial-block} with $a=\varpi_b$
on a separate block of $d-1$ variables. This gives rows
$g_1,\ldots,g_{3k}$ and weights $\tau_i>1/4$ with $\sum_i\tau_i=k$, and
$\sum_i\tau_ig_i$ is a sum of block forms that vanishes only at the product
$p_\star$ of the block singletons. Replace the rows by
$g'_i=g_i+\varepsilon'\sum_jg_j$ with $\varepsilon'=1/(8k)$. The sum of all
block Hessians is positive definite on $\R^n$, so each $g'_i$ has a positive
definite Hessian. The mixing matrix $I+\varepsilon'\mathbf1\mathbf1^{\mathsf
T}$ is invertible, so the span is unchanged, and it equals $3k$ because each
$T^d-\varpi_b$ is Eisenstein at $\varpi_b$ and the block spans have disjoint
supports. The weights $\eta_i=\tau_i-\varepsilon'k/(1+3k\varepsilon')=
\tau_i-1/11>1/8$ satisfy $\sum_i\eta_ig'_i=\sum_i\tau_ig_i$, so the mixed rows
still have the single feasible point $p_\star$. Each $g'_i$ is strictly convex
with rational coefficients, so its minimizer is rational and differs from the
irrational point $p_\star$, where $g'_i$ vanishes; hence each set
$\{g'_i\le0\}$ is a full-dimensional ellipsoid.

Next replace the first coordinate $x_{1,1}$ of the first block by
$y_1=x_{1,1}+\sum_{b\ge2}x_{b,1}+R_0$, keeping all other coordinates. The
linear part is unimodular, so positive definiteness, the span and the single
feasible point are preserved, and the new first coordinate of the feasible
point is $R_0+s$ with $s=\sum_b\varpi_b^{1/d}$. By Lemma~\ref{lem:qc-kummer},
$s$ has degree $d^k$. Every conjugate of $s$ has modulus at most
$\sum_b\varpi_b^{1/d}<R_0$, so every root $z_j$ of the minimal polynomial
$\prod_j(X-z_j)$ of $R_0+s$ has positive real part. Then
$\prod_j(X+z_j)$ is a product of linear factors $X+z_j$ with $z_j>0$ and
quadratic factors $X^2+2\operatorname{Re}(z)X+\abs z^2$ with
$\operatorname{Re}z>0$, so all its coefficients are positive. Since
$\prod_j(X-z_j)=(-1)^{d^k}\prod_j(-X+z_j)$, the coefficients of the minimal
polynomial alternate in sign and none vanishes.

Each block coefficient has $O(\log(d\varpi_k))$ bits. In $g'_i$ a nonconstant
coefficient involves only the three rows of one block, while a constant term
sums $3k$ constants; the shear changes each coefficient by a sum of at most
four products with $1$ and $R_0$. This gives the stated bit bounds, and there
are $3k$ rows with $O(n^2)$ coefficients each; in the explicit format each
monomial also carries an exponent vector of $O(n)=O(kd)$ bits. For the last
assertion, fix $k>3C$, let $\varpi_1,\ldots,\varpi_k$ be the first $k$ primes,
and let $d$ run through odd primes. Then $h=3k$, and $L=O_k(d^2\log d)$ in
the native encoding and $L=O_k(d^3)$ in the explicit format, while the
requested list for the first coordinate has $d^k+1$ entries, which eventually
exceeds $\phi(3k)L^C$ in both formats.
\end{proof}

In the unsheared system each coordinate has degree $d$, so its coordinatewise
output is short, while the common-field representation already needs degree
$d^k=(n/k+1)^k$. The shear moves this joint degree into one input coordinate.
The point also has a short description as a vector of real radicals, which
the corollary does not exclude. Unlike Theorem~\ref{thm:qc-sparse}, no
objective and no existential choice of weights or translations is involved.

\subsection{A single convex quartic row}
\label{app:qc-quartic}

The degree bound of Theorem~\ref{thm:qc-degree} depends on the span of the
constant Hessians of quadratic rows. For higher degree there are several
candidate parameters. The following example shows that the number of nonlinear
rows, or the dimension of the span of the polynomial Hessian matrices, does not
control the degree once a row has degree four. This is one reason why the quartic
singletons of the main text can have large degree with very few rows.

\begin{proposition}[One quartic row]
\label{prop:qc-quartic}
Let $\varpi_1,\ldots,\varpi_n$ be distinct primes, $\alpha_i=\varpi_i^{1/3}$,
$\Psi(x)=\sum_i(x_i^4-4\varpi_ix_i)$, and $M=3\sum_i\varpi_i^2$. The problem
\[
\min\{\tau: \Psi(x)-\tau\le0,\ 1\le x_i\le\varpi_i\ (1\le i\le n),\
-M\le\tau\le0\}
\]
has one globally convex nonlinear row, a strictly feasible point, and the
unique optimizer $x=\alpha$, $\tau=\theta=-3\sum_i\varpi_i\alpha_i$. The value
$\theta$ has degree exactly $3^n$. On the box, $\nabla^2\Psi\succeq12I$.
\end{proposition}

\begin{proof}
The identity
$x^4-4\varpi x+3\varpi\alpha=(x-\alpha)^2\bigl((x+\alpha)^2+2\alpha^2\bigr)$,
with $\alpha^3=\varpi$, shows that $\Psi$ has the unique minimizer $\alpha$, which
lies inside the box because $1<\alpha_i<\varpi_i$. Also $-M<\theta<0$, so the
bounds on $\tau$ are inactive, and $x_i=3/2$, $\tau=-1$ is strictly feasible:
$\Psi=\sum_i(81/16-6\varpi_i)\le-111n/16<-1$. The Hessian of the row is
$\diag(12x_1^2,\ldots,12x_n^2,0)\succeq0$.

Lemma~\ref{lem:qc-kummer} with $d=3$ and $c_i=-3\varpi_i$ shows that
$\theta$ has degree exactly $3^n$; it applies to every set of distinct primes,
including $3$.
\end{proof}

The polynomial Hessian of the single row spans a one-dimensional space of
polynomial matrices, while its coefficient matrices, or its values on the
box, span $n$ dimensions. The example refutes extensions of the quadratic
bound based on the number of rows or on the first notion, not on the second.
The exact convex quadratic lift $x_i^2\le y_i$,
$\sum_iy_i^2-4\sum_i\varpi_ix_i\le\tau$ has Hessian span $n+1$, in agreement
with Theorem~\ref{thm:qc-degree}. With the first $n$ primes and an indexed
sparse encoding, $L=\Theta(n\log n)$, so the degree is
$2^{\Theta(L/\log L)}$. A dense coefficient list of the quartic has
$\Theta(n^4)$ entries; the degree is then still superpolynomial in $L$, but
the exponential growth is in $n$, not in $L$. The value is also the short
radical sum displayed above,
and its sign relative to a rational threshold is a comparison problem that
the degree does not settle.

\subsection{Common-field representations from certified approximations}
\label{app:qc-recovery}

Degree and height bounds become an exact output only together with a way to
approximate one fixed point. The theorem below converts such data into a
common-field representation with polynomial overhead. It never multiplies
the degrees of individual coordinates, and it needs no factorization over a
number field. The only external tool is certified recognition of an algebraic
number from an approximation.

\begin{theorem}[Recognition of algebraic numbers~\cite{KannanLenstraLovasz1988}]
\label{thm:qc-kll}
There is a deterministic algorithm that, given integers $e,A\ge1$ and a
rational number $\tilde\vartheta$ with
$\abs{\tilde\vartheta-\vartheta}<2^{-\varsigma}/(12e)$ for a real algebraic
number $\vartheta$ whose primitive minimal polynomial has degree at most $e$
and coefficients bounded by $A$ in absolute value, where $\varsigma$ is the
least integer with $2^\varsigma>2^{e^2/2}(e+1)^{(3e+4)/2}A^{2e}$, returns that
minimal polynomial in time polynomial in $e$, $\log A$, and the length of
$\tilde\vartheta$.
\end{theorem}

\begin{theorem}[Common-field recovery]
\label{thm:qc-recovery}
Let $\beta\in\R^n$, $n\ge1$, and $\mathbb E=\Q(\beta_1,\ldots,\beta_n)$. Suppose
that bounds $\mathcal D\ge[\mathbb E:\Q]$ and $H\ge1$ are given, that every
$\beta_j$ has a primitive minimal polynomial with coefficients at most $2^H$
in absolute value, and that an oracle returns, for each integer $q\ge1$,
rationals $\tilde\beta_j$ with $\abs{\tilde\beta_j-\beta_j}<2^{-q}$ and length
polynomial in $n,\mathcal D,H,q$. A deterministic algorithm with polynomially
many oracle calls, at accuracies polynomial in $n,\mathcal D,H$, and with
additional time polynomial in $n,\mathcal D,H$, returns the actual degree
$d=[\mathbb E:\Q]$, a primitive irreducible integer polynomial $P$ of degree
$d$ with an isolating interval for a real root $\alpha$, integers $c_j$ with
$\alpha=\sum_jc_j\beta_j$, and rational polynomials $b_j$ of degree less than
$d$ with $\beta_j=b_j(\alpha)$. All output lengths are polynomial in
$n,\mathcal D,H$.
\end{theorem}

\begin{proof}
\emph{Heights.} Let $\mathfrak a_j$ be the absolute leading coefficient of the
minimal polynomial of $\beta_j$ and $\mathfrak A=\prod_j\mathfrak a_j$, so
$\log_2\mathfrak A\le nH$ and every $\mathfrak A\beta_j$ is an algebraic
integer. All conjugates of $\beta_j$ have modulus at most $2^{H+1}$. For
integers $\abs{c_j}\le2^{\mathfrak c}$ and $0\le\mathfrak s\le\mathcal D$, the
number $\vartheta=\sum_jc_j\beta_j+\mathfrak s\beta_i$ has conjugates of modulus at
most $\mathcal M=(n2^{\mathfrak c}+\mathcal D)2^{H+1}$, and
$\mathfrak A\vartheta$ is an algebraic integer. If $\mathfrak m_\vartheta$ is the monic
minimal polynomial of $\vartheta$, of degree $e$, then
$\mathfrak A^e\mathfrak m_\vartheta(T)$ is the monic minimal polynomial of
$\mathfrak A\vartheta$ evaluated at $\mathfrak AT$, so it has integer
coefficients bounded by $\mathfrak A^e(1+\mathcal M)^e$. Its primitive part is
the primitive minimal polynomial of $\vartheta$. Hence recognition can be
called with degree bound $\mathcal D$ and height
$\log_2A\le\mathcal D(nH+H+\mathfrak c+\log_2(n+\mathcal D)+3)$, using an
approximation of $\vartheta$ obtained from the oracle with $\log_2(n2^{\mathfrak
c}+\mathcal D)+1$ extra bits.

\emph{A primitive element.} Let $\alpha_\kappa=\sum_j\kappa^{j-1}\beta_j$ for
$0\le\kappa\le(n-1)\mathcal D(\mathcal D-1)/2$. For distinct embeddings
$\sigma,\tau$ of $\mathbb E$, the difference of their images of $\alpha_\kappa$
is a nonzero polynomial in $\kappa$ of degree at most $n-1$, so one candidate
separates all $d(d-1)/2$ pairs and generates $\mathbb E$. Recognize all
candidates; the largest degree found is $d$, and a candidate attaining it is a
primitive element $\alpha$ with its coefficients $c_j$. Its primitive minimal
polynomial $P$ has integer discriminant of modulus at least one, which gives a
root separation $\delta_P=2^{-4d^2(H_P+2)}$ when its coefficients are bounded
by $2^{H_P}$; an approximation of $\alpha$ within $\delta_P/8$ yields an
interval of half-length $\delta_P/4$ containing $\alpha$ and no other root.

\emph{Norm polynomials.} Fix $i$ and let $\gamma_{\mathfrak s}=\alpha+
\mathfrak s\beta_i$ for $\mathfrak s=0,\ldots,d$. Recognize the monic minimal
polynomial $\mathfrak m_{\mathfrak s}$ of $\gamma_{\mathfrak s}$, of degree
$e_{\mathfrak s}\mid d$, and put
$\Pi_{\mathfrak s}=\mathfrak m_{\mathfrak s}^{d/e_{\mathfrak s}}$. Each embedding of $\Q(\gamma_{\mathfrak s})$ extends to
exactly $d/e_{\mathfrak s}$ embeddings of $\mathbb E$, so
\[
\Pi_{\mathfrak s}(T)=\prod_{\sigma:\mathbb E\to\C}\bigl(T-\sigma(\gamma_{\mathfrak s})\bigr).
\]
Hence the $\Pi_{\mathfrak s}$ are the values at $S=\mathfrak s$ of
\[
\mathcal R_i(S,T)=\prod_{\sigma:\mathbb E\to\C}\bigl(T-\sigma(\alpha)-S\sigma(\beta_i)\bigr)
\in\Q[S,T],
\]
which is monic of degree $d$ in $T$ and has degree at most $d$ in $S$, and
coefficientwise Lagrange interpolation at $\mathfrak s=0,\ldots,d$ recovers it.
Interpolating the minimal polynomials themselves would be wrong when a sample
fails to generate $\mathbb E$: for $\mathbb E=\Q(\sqrt2)$, $\alpha=\sqrt2$ and
$\beta_i=-\sqrt2$, the sample at $\mathfrak s=1$ is $0$, with minimal polynomial
$T$ but norm polynomial $T^2$.

\emph{Coordinates.} Since $\alpha$ is primitive,
$\mathcal R_i(0,T)=\mathfrak m(T)$ is its monic minimal polynomial, and $\mathcal R_i(S,\alpha+S\beta_i)=0$ in
$\mathbb E[S]$ because the factor of the identity embedding vanishes.
Differentiating in $S$ at $S=0$ gives
\[
\partial_S\mathcal R_i(0,\alpha)+\mathfrak m'(\alpha)\beta_i=0,\qquad
\beta_i=-\frac{\partial_S\mathcal R_i(0,\alpha)}{\mathfrak m'(\alpha)} .
\]
Separability gives $\mathfrak m'(\alpha)\ne0$. Compute the inverse of
$\mathfrak m'$ modulo $\mathfrak m$ by solving the rational linear system of multiplication by $\mathfrak m'$ on
$\Q[T]/(\mathfrak m)$, and let $b_i$ be the remainder of
$-\mathfrak m'(T)^{-1}\partial_S\mathcal R_i(0,T)$ modulo $\mathfrak m$.

\emph{Cost.} There are $O(n\mathcal D^2)$ candidates and $n(d+1)$ samples.
Each recognized polynomial has degree at most $\mathcal D$ and coefficient
bits bounded as above. Raising to the power $d/e_{\mathfrak s}$, Lagrange
interpolation at integer nodes, whose denominators divide $\mathfrak
s!(d-\mathfrak s)!$, and the linear algebra modulo $\mathfrak m$ have polynomial bit
cost by determinant bounds.
\end{proof}

Reading coordinates from the derivative of a norm polynomial in a separating
parameter is the mechanism of rational univariate
representations~\cite{Rouillier1999}; the theorem combines it with certified
recognition, so that the input is an approximation oracle rather than a
polynomial system. The input must approximate one specified tuple. Separate minimal polynomials
of the coordinates, without selected embeddings, do not determine
$\mathbb E$, and a joint-degree bound alone does not give the representation.
The output lets a verifier evaluate any rational polynomial at $\beta$ by
arithmetic modulo $P$ and one sign determination at $\alpha$
(Lemma~\ref{lem:qc-sign}).

\begin{corollary}[Canonical feasible points at fixed span]
\label{cor:qc-feasible-point}
For each fixed $k$, a deterministic algorithm running in time polynomial in
$L$ decides whether a rational native PSD system with Hessian span $k$ is
feasible and, if so, returns its minimum-norm point in common-field format,
with field degree at most $\mathcal B(n,k)$ and all coefficients of
$L^{O(k+1)}$ bits.
\end{corollary}

\begin{proof}
Decide feasibility by Proposition~\ref{prop:qc-decision}, and let $p$ be the
minimum-norm point. Theorems~\ref{thm:qc-degree} and~\ref{thm:qc-height},
applied with $f=0$, give $\mathcal D=\mathcal B(n,k)$ and $H=L^{O(k+1)}$.
For the oracle, binary search the value $\norm p^2$ in $[0,nR_F^2]$ with
feasibility queries for $F\cap[-R_F,R_F]^n\cap\{\norm x^2\le\varsigma\}$,
whose Hessian span is at most $k+1$, until a feasible threshold
$\varsigma\le\norm p^2+2^{-2q-4}$ is found. Every point $x$ of this set satisfies
$\langle p,x-p\rangle\ge0$, so $\norm{x-p}^2\le\norm x^2-\norm p^2\le
2^{-2q-4}$. Bisecting coordinate intervals while keeping a nonempty
intersection with this set, again by feasibility queries, until all widths
are at most $2^{-q}$ gives a rational point within $2^{-q-1}+2^{-q-2}<2^{-q}$ of
$p$ in every coordinate. Each query has length polynomial in $L$ and $q$.
Theorem~\ref{thm:qc-recovery} finishes.
\end{proof}

For the canonical optimizer of a nonzero convex objective, the same
conversion applies to any certified approximation oracle for that optimizer.
Constructing such an oracle requires an exact optimization procedure for the
fixed-span class, which is part of the structural Hessian-span program and is
not developed here.

\subsection{One-field optimality certificates}
\label{app:qc-certificates}

An exact optimizer is most useful when a third party can check it without
repeating the optimization. When Slater's condition fails, the canonical
optimizer may have no Karush--Kuhn--Tucker multipliers for the original rows,
and it may be irrational. The certificate below handles both difficulties. It
supplies the optimizer in one real number field, adds a short sequence of
exposing records that preserve every feasible point, and ends with an exact
quadratic identity proving global optimality. All multipliers lie in the field
of the optimizer, so the field never grows, and at most $k$ curved records are
needed. The mechanisms (convex alternatives, facial reduction, polyhedral
normal cones) are classical; the point here is their encoding size in terms of
the Hessian span, and a verifier that needs only univariate arithmetic.

\begin{lemma}[Signs at an isolated root]
\label{lem:qc-sign}
Let $P\in\Z[T]$ have degree $e$ and coefficients at most $2^{H_P}$ in
absolute value, and let $\delta_P=2^{-4e^2(H_P+2)}$. Given rationals
$\varrho_-<\varrho_+$, one can check in polynomial time that $P$ is
squarefree, that $P(\varrho_-)P(\varrho_+)<0$, and that
$\varrho_+-\varrho_-<\delta_P$. If these checks pass, the interval contains
exactly one root $\alpha$ of $P$, and the sign of $\varsigma(\alpha)$ for a
given $\varsigma\in\Q[T]$ can be determined in time polynomial in the lengths
of $P$, the endpoints and $\varsigma$.
\end{lemma}

\begin{proof}
Squarefreeness is $\gcd(P,P')=1$. For squarefree $P$, the discriminant is
a nonzero integer equal to the leading coefficient to the power $2e-2$ times
the product of the squared differences of the roots, which have modulus at
most $2^{H_P+1}$. Bounding all but one factor shows that distinct roots differ
by more than $\delta_P$. A sign change on an interval shorter than $\delta_P$
therefore isolates exactly one root. For the sign, let $\mathfrak g=\gcd(P,\varsigma)$. Its roots are simple
roots of $P$, so $\varsigma(\alpha)=0$ exactly when
$\mathfrak g(\varrho_-)\mathfrak g(\varrho_+)<0$. Otherwise let
$P_1=P/\mathfrak g$, scaled to be primitive, and clear denominators of
$\varsigma$. Then $\operatorname{Res}(P_1,\varsigma)$ is a nonzero integer equal
to a power of the leading coefficient of $P_1$ times the product of
$\varsigma$ over the roots of $P_1$. Bounding the other factors by Cauchy's
bound gives an explicit lower bound $\epsilon_\varsigma$ on
$\abs{\varsigma(\alpha)}$, of polynomial bit length. Bisecting the interval
until its length times a Lipschitz bound for $\varsigma$ is below
$\epsilon_\varsigma$ determines the sign by evaluation at an endpoint. Gcds,
resultants, and these evaluations are polynomial-time rational
computations~\cite{vonzurGathenGerhard2013}.
\end{proof}

The lemma does not trust a claim that $P$ is irreducible, and it needs no
multivariate root isolation. Separate isolating intervals for unrelated
coordinates would not support this kind of joint evaluation.

\begin{definition}[One-field certificate]
\label{def:qc-certificate}
Consider a rational native PSD system with polyhedron $\{Ax\le b\}$ and
objective $f$. A \emph{one-field certificate} consists of a squarefree
$P\in\Z[T]$ with an interval isolating a root $\alpha$, rational polynomials
$b_1,\ldots,b_n$ defining $p=(b_j(\alpha))_j$, and records
$(\lambda^{(j)},\mu^{(j)})$ for $0\le j\le t$, all entries given as rational
polynomials in $\alpha$. The polyhedra $P_j=\{x:A_jx\le\beta_j\}$ are defined
by the verifier: $P_0$ consists of the original rows and the tangent rows
\[
g_i(p)+\nabla g_i(p)^{\mathsf T}(x-p)\le0\qquad(1\le i\le m),
\]
and $P_{j+1}$ adds to $P_j$ the equations $H_j(x-p)=0$ and $v_j^{\mathsf
T}(x-p)=0$, each written as two rows, where $H_j=\sum_i\lambda^{(j)}_iQ_i$ and
$v_j=\sum_i\lambda^{(j)}_i\nabla g_i(p)$. The verifier accepts if $p$
satisfies all original rows, if each exposing record $j<t$ satisfies
\begin{gather*}
\lambda^{(j)}\ge0,\quad\sum_i\lambda^{(j)}_i=1,\quad
\lambda^{(j)}_ig_i(p)=0,\quad\mu^{(j)}\ge0,\\
\mu^{(j)}_\iota(A_jp-\beta_j)_\iota=0,\quad
\sum_i\lambda^{(j)}_i\nabla g_i(p)+A_j^{\mathsf T}\mu^{(j)}=0,
\end{gather*}
and if the final record satisfies the same sign and complementarity
conditions without normalization, together with
$\nabla f(p)+\sum_i\lambda^{(t)}_i\nabla g_i(p)+A_t^{\mathsf T}\mu^{(t)}=0$.
\end{definition}

\begin{lemma}[Soundness]
\label{lem:qc-sound}
If the verifier accepts, then $p$ is feasible and globally optimal.
\end{lemma}

\begin{proof}
By convexity each $g_i$ lies above its tangent, so $F\subseteq P_0$. Suppose
$F\subseteq P_j$, and let $\mathcal G=\sum_i\lambda^{(j)}_ig_i$. For $x\in P_j$
the exact Taylor formula and the record give
\[
\mathcal G(x)=\tfrac12(x-p)^{\mathsf T}H_j(x-p)+
\sum_\iota\mu^{(j)}_\iota\bigl(\beta_j-A_jx\bigr)_\iota ,
\]
because $\mathcal G(p)=0$ by complementarity and $\nabla\mathcal G(p)=
-A_j^{\mathsf T}\mu^{(j)}$ with $\mu^{(j)}$ supported on rows tight at $p$. Both
terms are nonnegative on $P_j$, and $\mathcal G\le0$ on $F$. Hence both vanish
on $F$; positive semidefiniteness turns the first into $H_j(x-p)=0$, and the
second gives $v_j^{\mathsf T}(x-p)=0$. Thus $F\subseteq P_{j+1}$. For the final
record and $x\in F\subseteq P_t$,
\[
f(x)-f(p)=\tfrac12(x-p)^{\mathsf T}\Bigl(Q_0+\sum_i\lambda^{(t)}_iQ_i\Bigr)(x-p)
-\sum_i\lambda^{(t)}_ig_i(x)+\sum_\iota\mu^{(t)}_\iota(\beta_t-A_tx)_\iota\ge0 .
\qedhere
\]
\end{proof}

The verifier checks neither a constraint qualification, nor minimality of the
norm of $p$, nor that a record reduces a dimension. Complementarity with the
quadratic rows is essential: without it, $\mathcal G(p)$ need not vanish. The
sets $P_{j+1}$ are polyhedral sections containing $F$ and need not be faces of
$P_j$: the record $\mathcal G=x^2$ on $P_0=\R$ yields $\{0\}$.

\begin{theorem}[One-field certificates]
\label{thm:qc-certificate}
Let a rational native PSD system with Hessian span $k$ and a convex quadratic
objective attain a finite optimum, and let $p$ be the canonical optimizer.
There is a one-field certificate for $p$ in the field $\Q(p)$, of degree at
most $\mathcal B(n,k)$, with at most $k$ exposing records, each having at most
$n+1$ positive entries in $(\lambda^{(j)},\mu^{(j)})$, and a final record with at
most $n$ positive entries. Its total length is $L^{O(k+1)}$, and the verifier
runs in time polynomial in $L$ and the certificate length.
\end{theorem}

\begin{proof}
\emph{Existence of the records.} Let $\mathcal V_j$ be the direction space of
the affine hull of $P_j$, and let $I_j$ contain the quadratic rows whose
Hessian restricted to $\mathcal V_j$ is nonzero. A row outside $I_j$
coincides on $P_j$ with its tangent row, so it is implied by $P_j$, and
$F=\{x\in P_j:g_i(x)\le0,\ i\in I_j\}$.

If the rows in $I_j$ have a common point of $P_j$ at which all of them are
negative, including the case $I_j=\varnothing$, the convex
Karush--Kuhn--Tucker theorem with affine constraints and a strict point for
the nonlinear ones~\cite{Rockafellar1970} gives multipliers at the optimizer
$p$, which form the final record. Otherwise the convex alternative, as in
Proposition~\ref{prop:qc-two-rows}, gives normalized weights $\lambda\ge0$ on
$I_j$ with $\mathcal G=\sum_i\lambda_ig_i\ge0$ on $P_j$. Since $p\in F$,
$\mathcal G(p)=0$, so positive weights sit on rows active at $p$, and $p$
minimizes $\mathcal G$ over the polyhedron $P_j$. Polyhedral first-order
optimality supplies $\mu\ge0$ on the rows of $P_j$ tight at $p$ with
$\nabla\mathcal G(p)+A_j^{\mathsf T}\mu=0$.

In both cases the multipliers solve a linear system with coefficients in
$\Q(p)$: $n$ stationarity equations, plus normalization for an exposing
record, in nonnegative unknowns indexed by the active rows. A nonnegative
solution of minimal support has linearly independent columns, so it has at
most $n+1$, respectively $n$, positive entries and is given by Cramer's rule
over $\Q(p)$.

For an exposing record, $H_j$ restricted to $\mathcal V_j$ is a positive
combination of nonzero positive semidefinite forms, hence nonzero, while
$\mathcal V_{j+1}\subseteq\ker H_j$. The restriction from $\mathcal V_j$ to
$\mathcal V_{j+1}$ therefore kills a nonzero element of the span of the
restricted Hessians, and the dimension of that span drops by at least one.
It starts at most at $k$, so there are at most $k$ exposing records.

\emph{Size.} By Theorems~\ref{thm:qc-degree} and~\ref{thm:qc-height},
$K=\Q(p)$ has degree $d\le\mathcal B(n,k)\le(2L)^k$ and the coordinates of $p$
have heights at most $B_*=L^{O(k+1)}$. The tangent rows have heights
$O(LB_*)$. In each round the new multipliers are Cramer quotients of minors of
order at most $n+1$, and the new rows are sums of polynomially many products
of these with earlier data. By Lemma~\ref{lem:qc-heights}, if all earlier data
have heights at most $B_j$, the new ones have heights at most
$cL^c(B_j+B_*)$ for an absolute $c$. After at most $k+1$ rounds all heights
are $L^{O(k)}B_*=L^{O(k+1)}$. Heights, unlike power-basis coefficients, do not
grow by a factor of the field degree in each round.

\emph{Encoding.} Choose $\alpha=\sum_jc_jp_j$ primitive with integers
$\abs{c_j}\le(1+(n-1)d^2)^{n-1}$, as in the proof of
Theorem~\ref{thm:qc-recovery}. Theorem~\ref{thm:qc-height}(i) and
\eqref{eq:qc-minpoly-height} bound its minimal polynomial $P$ by
$L^{O(k+1)}$ bits. For $\xi\in K$, write $\xi=\sum_{r<d}\varkappa_r\alpha^r$.
Cramer's rule for the Vandermonde system over the $d$ embeddings of $K$, with
Lemma~\ref{lem:qc-heights}(a),(d), gives
$\height(\varkappa_r)\le2d^3\height(\alpha)+d\height(\xi)+2\log d!$; the normal
closure enters only through invariance of heights under conjugation. Each
rational $\varkappa_r$ therefore has $L^{O(k+1)}$ bits. There are polynomially
many entries in the records, so the total length is $L^{O(k+1)}$.

\emph{Verification.} The verifier forms the tangent rows and the sections by
rational polynomial arithmetic modulo $P$, checks the identities by reduction
modulo $P$, and checks every sign with Lemma~\ref{lem:qc-sign}.
\end{proof}

Up to $k$ exposing records may be needed. For the chain
$x_i^2-x_{i+1}\le0$ ($1\le i<k$), $x_k^2\le0$, with objective $x_1$, the
feasible set is $\{0\}$, and the original rows admit no multipliers at $0$.
The tangent rows give $x_2,\ldots,x_k\ge0$. If an exposing record put positive
weight $\lambda_i$ on a row $i<k$ whose variable $x_{i+1}$ is not yet fixed to
zero, then along $\tau e_{i+1}$ the aggregate would equal
$-\lambda_i\tau+O(\tau^2)<0$ for small $\tau>0$. Hence each record can fix
only the last free variable, and $k$ records are needed before the final
identity can certify the objective $x_1$.

The theorem is an existence and verification statement. The primal point can
be constructed by Corollary~\ref{cor:qc-feasible-point} when $f=0$, but no
polynomial-time construction of the exposing multipliers is claimed. A
certificate applied after fixing an integer assignment certifies optimality
in that continuous fiber only, not global mixed-integer optimality. For an
empty feasible set, the next subsection gives a rational certificate.

\subsection{Rational infeasibility certificates}
\label{app:qc-infeasible}

An infeasible native PSD system has a positive aggregate: a convex
combination of its rows that is bounded below by a positive constant. This
alternative is classical, without any constraint qualification or
compactness~\cite{JeyakumarLi2014}. Proposition~\ref{prop:qc-span-three}
shows that aggregates with minimum \emph{zero} may need irrational weights.
With a \emph{positive} margin, the weights can be chosen rational, and their
length is controlled by the Hessian span. Such a certificate is checked with
rational arithmetic alone.

\begin{theorem}[Rational infeasibility certificates]
\label{thm:qc-rational-infeasibility}
Let rows $g_i(x)=\tfrac12x^{\mathsf T}Q_ix+a_i^{\mathsf T}x+c_i\le0$,
$1\le i\le m$, with rational data and $Q_i\succeq0$, have no common solution
in $\R^n$; affine rows are included with $Q_i=0$, and equalities as two rows.
Let $k$ be the Hessian span and $L$ the input length. There are rational
weights $w_i\ge0$ with $\sum_iw_i=1$, at most $n+1$ of them positive, a
rational $z\in\Q^n$ and a rational $\gamma>0$ such that
\begin{equation}
\label{eq:qc-rational-aggregate}
\sum_iw_ig_i(x)=\gamma+\tfrac12(x-z)^{\mathsf T}H(x-z),\qquad
H=\sum_iw_iQ_i\succeq0,
\end{equation}
of total length $L^{O(k+1)}$. A verifier checks $w\ge0$, $\sum_iw_i=1$,
$Hz=-\sum_iw_ia_i$ and $\gamma=\sum_iw_ic_i+\tfrac12(\sum_iw_ia_i)^{\mathsf
T}z>0$ in polynomial time.
\end{theorem}

\begin{proof}
\emph{An algebraic aggregate.} The epigraph problem
$\min\{\tau:g_i(x)\le\tau\ (1\le i\le m)\}$ is feasible, and its infimum is
nonnegative because the system is infeasible. A feasible convex
quadratically constrained problem with a linear objective and a finite
infimum attains it, without compactness or a Slater
condition~\cite[Corollary~2, in the numbering of the 1997
report]{LuoZhang1999}; the attained value $\alpha$ is positive, since
$\alpha=0$ would give a feasible point. Let $(\bar x,\alpha)$ be the canonical
optimizer. The epigraph rows have Hessians $\diag(Q_i,0)$, so
Theorems~\ref{thm:qc-degree} and~\ref{thm:qc-height}, in dimension $n+1$, show
that $K'=\Q(\bar x,\alpha)$ has degree at most $\mathcal B(n+1,k)$ and that
$\bar x$ and $\alpha$ have heights $L^{O(k+1)}$. The epigraph is strictly
feasible, so the convex Karush--Kuhn--Tucker conditions give weights
$w_i^*\ge0$ with $\sum_iw_i^*=1$, $\sum_iw_i^*\nabla g_i(\bar x)=0$, and
$w_i^*=0$ unless $g_i(\bar x)=\alpha$. A basic solution of this system of
$n+1$ equations has at most $n+1$ positive entries, lies in $K'$, and has
heights $L^{O(k+1)}$ by Cramer's rule. The aggregate
$\mathcal G^*=\sum_iw_i^*g_i$ is stationary at $\bar x$ with value $\alpha$, so
$\mathcal G^*\ge\alpha$. Lemma~\ref{lem:qc-heights}(c) and the degree bound
give
\[
\omega=\min\{w_i^*:w_i^*>0\}\ge2^{-L^{O(k+1)}},\qquad\alpha\ge2^{-L^{O(k+1)}} .
\]

\emph{Rounding inside the exact equations.} Let $I$ be the positive support
and $\Omega=\bigcap_{i\in I}\ker Q_i$, a rational subspace. For every weight
vector that is positive on $I$, the kernel of $\sum_{i\in I}w_iQ_i$ equals
$\Omega$, because a sum of positive semidefinite forms vanishes only where
each summand does. Such an aggregate has a finite minimum exactly when
$\sum_iw_ia_i\perp\Omega$. Hence we round inside the rational affine space
\[
\mathcal W=\Bigl\{w\in\R^I:\sum_iw_i=1,\ \sum_iw_ia_i\perp\Omega\Bigr\},
\]
which contains $w^*$. Independent rounding of the weights would not do: for
$x^2+y+1$ and $x^2-y+1$, equal weights give $x^2+1$, while weights
$(\tfrac12+\epsilon,\tfrac12-\epsilon)$ give an aggregate unbounded below.
Parametrize $\mathcal W$ rationally with free coordinates chosen among the
weights, round the free coordinates to a dyadic grid, and let $\hat w$ be the
result; then $\varepsilon=\norm{\hat w-w^*}_1$ can be made as small as
$2^{-\mathfrak b}$ with weights of $\mathfrak b+L^{O(1)}$ bits.

\emph{The margin survives.} Let $M\ge1$ bound $\norm{Q_i}$, $\norm{a_i}$ and
$\abs{c_i}$. If $\sum_{i\in I}Q_i\ne0$, let $\varrho$ be its least positive
eigenvalue. The product of the positive eigenvalues of an integer multiple of
this matrix is a positive integer coefficient of its characteristic
polynomial, and all eigenvalues are at most $2^{L^{O(1)}}$, so
$\varrho\ge2^{-L^{O(1)}}$. Put $\eta=\min(1,\omega\varrho)$. On
$\mathcal V_0=\Omega^\perp$ the aggregate Hessian at $w^*$ is at least $\eta
I$, and if $\varepsilon\le\omega/2$, at least $(\eta/2)I$ at $\hat w$, with the
same support and kernel. The minimum of the aggregate is
$\nu(w)=c(w)-\tfrac12b(w)^{\mathsf T}(H(w)|_{\mathcal V_0})^{-1}b(w)$, with
$b(w)=\sum_iw_ia_i$ and $c(w)=\sum_iw_ic_i$, and the identity
$\hat H^{-1}-H^{*-1}=\hat H^{-1}(H^*-\hat H)H^{*-1}$ gives
\[
\abs{\nu(\hat w)-\nu(w^*)}\le\Bigl(M+\frac{2M^2}{\eta}+\frac{M^3}{\eta^2}\Bigr)
\varepsilon\le\frac{4M^3}{\eta^2}\varepsilon .
\]
Choosing $\varepsilon\le\min\{\omega/4,\alpha\eta^2/(16M^3)\}$ keeps the
support and gives $\nu(\hat w)\ge3\alpha/4$. If $\sum_{i\in I}Q_i=0$, then
$\mathcal W$ forces $b(w)=0$, the aggregate is the constant $c(w)$, and
$\varepsilon\le\min\{\omega/4,\alpha/(4M)\}$ suffices. In both cases
$\mathfrak b=L^{O(k+1)}$.

\emph{The certificate.} The rational system $\hat Hz=-b(\hat w)$ is
consistent because $b(\hat w)\perp\Omega=\ker\hat H$. A rational solution
from Gaussian elimination and $\gamma=c(\hat w)+\tfrac12b(\hat w)^{\mathsf
T}z=\nu(\hat w)>0$ satisfy~\eqref{eq:qc-rational-aggregate}, and all lengths
are $L^{O(k+1)}$. A feasible point would make the left side
of~\eqref{eq:qc-rational-aggregate} nonpositive and the right side positive.
\end{proof}

The rounding step is an algorithm once an algebraic aggregate is known: its
support, the rational kernel, the precision, and the rounding are all
computable. Finding the initial aggregate is a separate task that the theorem
does not address. The certificate refutes real feasibility only; it says
nothing about an integer program whose continuous relaxation is feasible.

The span must enter exponentially for this format, and the support cannot be
bounded in terms of $k$ alone.

\begin{proposition}[Lower bounds for positive aggregates]
\label{prop:qc-infeasibility-lower}
\begin{enumerate}[label=(\alph*),leftmargin=*]
\item The system $\tfrac12-x_1\le0$, $x_i^2-x_{i+1}\le0$ for $1\le i<k$,
$x_k^2\le0$ is infeasible, has Hessian span $k$, and every normalized
positive aggregate has margin at most $2^{-2^k}$, so a rational margin has a
denominator of at least $2^k+1$ bits.
\item In the system $\tfrac12-x_1\le0$, $x_i^2-x_{i+1}\le0$ for $1\le i\le k$,
$x_{k+1}\le0$, with span $k$, every positive aggregate with rational weights
$\lambda_i\ge0$ satisfies $\lambda_{k+1}/\lambda_0>2^{d-1}/(3d)$, where
$d=2^k$; the two weights together need $d-O(\log d)$ bits under any common
scaling.
\item For $n\ge2$, the rows $\norm{x-e_i}^2-\frac{n-2}{n-1}\le0$,
$1\le i\le n$, have $k=1$, are infeasible, and every proper subfamily is
feasible, so every positive aggregate uses all $n$ rows.
\end{enumerate}
\end{proposition}

\begin{proof}
(a) Backward induction from $x_k=0$ gives $x_1=0$, contradicting the first
row. At $y_i=2^{-2^{i-1}}$ every row except the last vanishes and the last
equals $2^{-2^k}$; a normalized aggregate with margin $\gamma$ satisfies
$\gamma\le\sum_iw_ig_i(y)\le2^{-2^k}$.

(b) Infeasibility follows as in (a). At $x=0$ the aggregate equals
$\lambda_0/2$, so $\lambda_0>0$. On the curve $x_i=\tau^{2^{i-1}}$ all middle
rows vanish, and positivity at $\tau=(1+1/d)/2$ gives
$\lambda_0\bigl(-1/(2d)\bigr)+\lambda_{k+1}2^{-d}(1+1/d)^d>0$. Since
$(1+1/d)^d<3$, the ratio bound follows. If $\lambda_0=a_0/b_0$ and
$\lambda_{k+1}=a_1/b_1$ in lowest terms, then $a_1b_0>2^{d-1}/(3d)$.

(c) The average of the rows is $\norm{x-\mathbf1/n}^2+\frac{n-1}n-\frac{n-2}{n-1}
>0$. For a subfamily of $s\le n-1$ indices, the average of its centers has
squared distance $1-1/s\le\frac{n-2}{n-1}$ to each of them.
\end{proof}

These are lower bounds for one certificate format. Both chains have short
symbolic proofs of infeasibility, so no lower bound for other proof systems
follows.

\subsection{Number-field affine data}
\label{app:qc-number-field}

Theorem~\ref{thm:qc-height} allows coefficients in a number field. The
following theorem shows that a convex quadratic program whose matrices are
rational, but whose right-hand sides and linear costs lie in one explicitly
given real number field, is solved exactly inside that field in polynomial
time. The field degree is part of the input length; no product of degrees of
independently given coefficients is allowed. The argument uses only exact
rational quadratic programming, a determinant bound in the field, and two
error bounds.

\begin{theorem}[Rational matrices, algebraic affine data]
\label{thm:qc-number-field-qp}
Let $\mathbb F=\Q(\alpha)\subset\R$ be given by a dense irreducible integer
polynomial of degree $D$ and an interval isolating $\alpha$. Let $Q\succeq0$
and $C$ be rational, let $b\in\mathbb F^s$ and $c\in\mathbb F^n$ be given in the
power basis, and let $L$ be the total input length. A deterministic algorithm
running in time polynomial in $L$ decides whether
$\min\{\tfrac12x^{\mathsf T}Qx+c^{\mathsf T}x:Cx\le b\}$ is infeasible,
unbounded, or finite, and in the last case returns the minimum-norm optimizer
$p\in\mathbb F^n$ in the power basis, with polynomial length.
\end{theorem}

\begin{proof}
\emph{Charts.} Assume first that the problem is feasible and bounded below; a
convex quadratic bounded below on a nonempty polyhedron attains its
minimum~\cite{FrankWolfe1956}; this is also the polyhedral case
of~\cite[Corollary~2, in the numbering of the 1997 report]{LuoZhang1999}.
Let $C_I$ consist of rows forming a basis of all constraint normals active
at $p$, and let
$\mathcal M_I=\left(\begin{smallmatrix}Q&C_I^{\mathsf T}\\C_I&0\end{smallmatrix}
\right)$. For $\xi\in\ker C_I$, the points $p\pm\tau\xi$ are feasible for small
$\tau$, so $Qp+c\in(\ker C_I)^\perp$, and some signed $\lambda$ satisfies
$\mathcal M_I(p,\lambda)=(-c,b_I)$. The kernel of $\mathcal M_I$ is
$(\ker Q\cap\ker C_I)\times\{0\}$: if $Q\xi+C_I^{\mathsf T}\zeta=0$ and
$C_I\xi=0$, then $\xi^{\mathsf T}Q\xi=0$, so $Q\xi=0$ and $\zeta=0$. For such
$\xi$, $f(p\pm\tau\xi)=f(p)$, and minimality of the norm gives $p\perp\xi$.
Hence $(p,\lambda)$ is the least-norm solution, and
\[
(p,\lambda)=\mathcal M_I^+(-c,b_I),\qquad
\mathcal M_I^+=\mathcal A^{-1}\mathcal M_I\mathcal A^{-1},\quad
\mathcal A=\mathcal M_I+ZZ^{\mathsf T},
\]
where the columns of $Z$ form a rational basis of $\ker\mathcal M_I$; on the
kernel $\mathcal A$ acts invertibly and $\mathcal M_I$ vanishes, and on the
range $\mathcal A$ agrees with $\mathcal M_I$. Thus each coordinate of $p$ is a
rational linear combination of the entries of $b$ and $c$ with coefficients of
polynomial bit length, uniformly over all $I$. This proves $p\in\mathbb F^n$,
bounds its power-basis length, and gives a computable rational $R\ge1$ of
polynomial length with $p\in[-R,R]^n$, using Cauchy's bound for $\alpha$. The
chart is used only in the proof; the algorithm does not search over $I$.

\emph{Separation in $\mathbb F$.} Clearing denominators and multiplying by a
power of the leading coefficient of the polynomial of $\alpha$ makes every
entry of a given matrix over $\mathbb F$ an algebraic integer. A nonzero
$r\times r$ minor is then, up to a known integer factor, a nonzero algebraic
integer whose conjugates are bounded by determinant expansion, so its norm is
a nonzero integer. Hence every nonzero minor, and every nonzero element of
$\mathbb F$ of polynomial length, has absolute value at least $2^{-L^{O(1)}}$.

\emph{Error bounds.} For a real matrix $\mathfrak M$ and any right-hand side
$\mathfrak r$ with $\{\mathfrak Mx\le\mathfrak r\}\ne\varnothing$,
\begin{equation}
\label{eq:qc-hoffman}
\dist\bigl(y,\{\mathfrak Mx\le\mathfrak r\}\bigr)\le
H_{\mathfrak M}\norm{(\mathfrak My-\mathfrak r)_+}_\infty,\qquad
H_{\mathfrak M}=\max_{\mathcal I}\frac{\sqrt{\abs{\mathcal I}}}
{\sigma_{\min}(\mathfrak M_{\mathcal I})},
\end{equation}
the maximum over linearly independent sets of rows~\cite{Hoffman1952}. Indeed,
if $z$ is the projection of $y$, then $y-z=\mathfrak M_{\mathcal I}^{\mathsf
T}\lambda$ with $\lambda\ge0$ and independent rows active at $z$, and
$\norm{y-z}^2=\lambda^{\mathsf T}(\mathfrak M_{\mathcal I}y-\mathfrak
r_{\mathcal I})\le\norm\lambda_1\norm{(\mathfrak My-\mathfrak r)_+}_\infty$,
with $\norm\lambda_1\le\sqrt{\abs{\mathcal I}}\,\norm{y-z}/\sigma_{\min}$. A
nonsingular square column submatrix $\mathfrak N$ of $\mathfrak M_{\mathcal
I}$ gives $\sigma_{\min}(\mathfrak M_{\mathcal I})\ge\abs{\det\mathfrak N}/
\norm{\mathfrak N}^{\abs{\mathcal I}-1}$. The separation bound for nonzero
minors and entrywise bounds for $\norm{\mathfrak N}$ hold uniformly over all
row sets, so they yield rational numbers $H_A,H_{\mathfrak M}\ge1$ of
polynomial length that bound the constant in~\eqref{eq:qc-hoffman} for the
rational matrix $A$ of the rows $Cx\le b$ and $\norm x_\infty\le R$, and for
$\mathfrak M=(A;Q;-Q;c^{\mathsf T};-c^{\mathsf T})$;
inequality~\eqref{eq:qc-hoffman} remains true with these upper bounds. The
right-hand sides do not enter.

Let $P_R=\{Cx\le b,\ \norm x_\infty\le R\}$, which has the same optimum and
canonical optimizer, let $\mathcal O$ be its optimal set, and let $U=nR$, a
rational bound for $\norm x$ on the box. Two optimal points have the same
$Qx$, by the midpoint identity, and then the same $c^{\mathsf T}x$;
conversely such points are optimal. So $\mathcal O=\{x\in P_R:Qx=Qp,\ c^{\mathsf T}x=c^{\mathsf T}p\}$,
whose matrix is $\mathfrak M$. For $x\in P_R$ with $\gamma=f(x)-f(p)\le1$, the
variational inequality at $p$ gives $\tfrac12(x-p)^{\mathsf T}Q(x-p)\le\gamma$,
hence $\norm{Q(x-p)}\le\sqrt{2\mathfrak q\gamma}$ for a rational
$\mathfrak q\ge\norm Q$, and
$\abs{c^{\mathsf T}(x-p)}\le\gamma+U\sqrt{2\mathfrak q\gamma}$. By
\eqref{eq:qc-hoffman},
\[
\dist(x,\mathcal O)\le\mathcal T\sqrt\gamma,\qquad
\mathcal T=H_{\mathfrak M}\bigl(1+2(U+1)(\mathfrak q+1)\bigr).
\]

\emph{Approximating $p$.} For rational $\varepsilon,\delta>0$, compute
rational $b^+$ with $b\le b^+\le b+\varepsilon\mathbf1$ and $\tilde c$ with
$\norm{\tilde c-c}\le\varepsilon$ by exact root refinement of $\alpha$, keep
the box exact, and solve the strictly convex rational problem
\[
y=\argmin\bigl\{\tfrac12x^{\mathsf T}Qx+\tilde c^{\mathsf T}x+
\tfrac\delta2\norm x^2:\ Cx\le b^+,\ \norm x_\infty\le R\bigr\}
\]
exactly in polynomial time~\cite{KozlovTarasovKhachiyan1980}. Let
$L_f=\mathfrak qU+\bar c$, where $\bar c\ge1$ is a rational upper bound for
$\norm c_1$ obtained from the power-basis coefficients of $c$ and Cauchy's
bound for $\alpha$; it is a Lipschitz bound for $f$ on the box. Put
$\mathcal W=2U+L_fH_A$. The constants $U$, $\mathfrak q$, $L_f$, $H_A$,
$H_{\mathfrak M}$, $\mathcal T$ and $\mathcal W$ are rationals of polynomial
length, computed before $\delta$ and $\varepsilon$ are chosen below. Some
$x\in P_R$ satisfies $\norm{x-y}\le H_A\varepsilon$, and comparison with $p$
gives
\[
\norm y^2\le\norm p^2+2\mathcal W\varepsilon/\delta,\qquad
f(x)-f(p)\le\Gamma:=\mathcal W\varepsilon+\delta U^2/2 .
\]
If $\Gamma\le1$, then $e:=\dist(y,\mathcal O)\le H_A\varepsilon+\mathcal T
\sqrt\Gamma$. Let $\hat y$ be the projection of $y$ onto $\mathcal O$. Since
$p$ is the projection of $0$ onto $\mathcal O$,
$\norm{\hat y-p}^2\le\norm{\hat y}^2-\norm p^2\le2\mathcal W\varepsilon/\delta
+2Ue$, so
\[
\norm{y-p}\le e+\sqrt{2\mathcal W\varepsilon/\delta+2Ue} .
\]
For a target accuracy $\tau=2^{-\mathfrak q'}\le1$ put
$\mathcal E=\tau^2/(16(U+1))$, $\delta=\mathcal E^2/(8\mathcal T^2(U^2+1))$ and
$\varepsilon\le\min\{\mathcal E/(2H_A),\mathcal E^2/(16\mathcal W\mathcal T^2),
\delta\tau^2/(32\mathcal W)\}$. Then $\Gamma\le\mathcal E^2/(8\mathcal T^2)$,
$e\le\mathcal E$, and $\norm{y-p}\le\tau/16+\sqrt3\tau/4<\tau$, with all
parameters of $O(\mathfrak q')+L^{O(1)}$ bits. The norm penalty is necessary:
without it, the problems $\min\{-\epsilon x:0\le x\le1\}$ select $1$ for
every $\epsilon>0$, while the canonical optimizer of the limit problem is $0$.

\emph{Exact output.} Every row slack $b_\iota-C_\iota p$ lies in $\mathbb F$
with polynomial length, so it is either zero or at least $\mathfrak g$ in
absolute value, where $\mathfrak g=2^{-L^{O(1)}}$ is a rational lower bound
computed from the separation bound. Put
$\mathfrak C=\max\{1,\max_\iota\norm{C_\iota}_1\}$, where the inner maximum is
$0$ if $C$ has no rows. Approximate $p$ within $\mathfrak g/(8\mathfrak C)$ in
the maximum norm and the entries of $b$ within $\mathfrak g/8$. Each computed
slack is then within $\mathfrak g/4$ of the true one, so a row is active
exactly when its computed slack is below $\mathfrak g/2$. Choose a maximal
independent subset $I$ of the active rows by rational linear algebra and
return the first block of $\mathcal M_I^+(-c,b_I)$, computed by rational
matrix arithmetic applied to power-basis vectors.

\emph{Classification.} For a rational matrix $\mathfrak B$ and
$\mathfrak a\in\mathbb F^{s'}$, feasibility of $\mathfrak Bx\le\mathfrak a$ is
decided as follows. The chart argument for $\norm x^2$ gives a box
$[-R',R']^{n}$ of polynomial length containing the least-norm feasible point
if one exists. The value $t^*=\min\{t\ge0:\mathfrak Bx\le\mathfrak a+t\mathbf1,
\ \norm x_\infty\le R'\}$ is attained at a vertex, so it lies in $\mathbb F$
with polynomial length; it is zero exactly when the system is feasible, and
otherwise at least $2^{-L^{O(1)}}$. Rounding $\mathfrak a$ outward by less
than this gap and solving one rational linear program decides
feasibility~\cite{GroetschelLovaszSchrijver1988}. For a feasible problem,
$f$ is bounded below exactly when $Qv+C^{\mathsf T}\lambda=-c$, $\lambda\ge0$,
is feasible: completing the square gives a lower bound, and otherwise
separation from the closed cone $Q\R^n+C^{\mathsf T}\R^s_+$ gives $h$ with
$Qh=0$, $Ch\le0$ and $c^{\mathsf T}h<0$, an unbounded ray. Both tests use the
same routine.
\end{proof}

The input model is one explicitly represented field. Coefficients given by
separate minimal polynomials could generate a compositum of exponential
degree, and that case is not covered. The existence of rational charts alone
does not give the algorithm; the approximation with a norm penalty does.

\subsection{A nonconvex integer caveat}
\label{app:qc-pell}

The bounds above use convexity of the native rows. Without it, a classical
Pell family shows that two integer variables and no continuous variable
already force exponentially long integer output. The phenomenon is
classical~\cite{Lenstra2002}; it is recorded only to mark why convexity
appears in the statements of this paper.

\begin{proposition}[Pell output]
\label{prop:qc-pell}
For $m\ge1$, the problem
\[
\min\{x:\ x^2-5^{2m+1}y^2=1,\ x\ge2,\ y\ge1,\ x,y\in\Z\}
\]
has input length $\Theta(m)$, no continuous variables, and the unique
optimizer $(X_{5^m},B_{5^m}/5^m)$, where $(9+4\sqrt5)^j=X_j+B_j\sqrt5$. Every
feasible point has an $x$-coordinate of $\Theta(5^m)$ or more bits, and the
optimal value has $\Theta(5^m)$ bits.
\end{proposition}

\begin{proof}
With $B=5^my$, the equation becomes $X^2-5B^2=1$. Every solution in positive
integers is $(X_j,B_j)$ for a unique $j\ge1$: dividing $X+B\sqrt5$ by a power of
$u=9+4\sqrt5$ brings it into $[1,u)$, where the quotient $U+V\sqrt5$ has
integer coefficients, norm one, $U\ge1$ and $0\le V<4$, and $1+5V^2$ is a
square only for $V=0$. The binomial expansion gives
$B_j\equiv4j9^{j-1}\pmod5$, and expanding the fifth power gives
$B_{5j}=5B_j(X_j^4+10X_j^2B_j^2+5B_j^4)$, whose second factor is congruent to
one modulo five because $X_j^2\equiv1$. Hence the five-adic valuation of
$B_j$ equals that of $j$. Feasibility requires $5^m\mid B_j$, that is,
$5^m\mid j$, and $X_j=(u^j+u^{-j})/2$ increases with $j$. Finally
$u^j/2<X_j<u^j$ with $16<u<32$.
\end{proof}

The optimal value is an integer, so its minimal polynomial is linear; the
obstruction is coefficient length, not degree. Repeated fifth powering gives
an arithmetic circuit with $O(m)$ gates for the optimizer, so the obstruction
concerns expanded output only. Each instance is feasible, and no hardness of
deciding feasibility follows. The rows are indefinite and the feasible set is
nonconvex, so the example does not bear on the convex mixed-integer results
of Section~\ref{sec:constraints}.
